% !TeX program = pdflatex
% Kandydacki skrypt; nie jest jeszcze zarejestrowany w Atlasie.
\documentclass[11pt,a4paper]{article}
\usepackage[T1]{fontenc}
\usepackage[utf8]{inputenc}
\usepackage[polish]{babel}
\usepackage{lmodern}
\usepackage{amsmath,amssymb,amsthm,mathtools}
\usepackage{microtype}
\usepackage{geometry}
\usepackage{enumitem}
\usepackage{tocloft}
\usepackage{xcolor}
\usepackage{tikz}
\usetikzlibrary{arrows.meta}
\usepackage{xr-hyper}
\usepackage[unicode]{hyperref}
\geometry{margin=2.5cm}
\hypersetup{colorlinks=true,linkcolor=blue!45!black,urlcolor=blue!45!black,pdftitle={Geometria różniczkowa w przestrzeniach funkcyjnych}}
\setlength{\parindent}{0pt}
\setlength{\parskip}{4pt}
\setlength{\emergencystretch}{2em}
\setlength{\cftsecnumwidth}{3em}
\setlength{\cftsubsecnumwidth}{4.5em}
% Moduły 15A/15B mają numery podpodsekcji w rodzaju 15B.12.8.
\setlength{\cftsubsubsecnumwidth}{6em}
\setlist[enumerate]{topsep=4pt,itemsep=2pt}
\setlist[itemize]{topsep=4pt,itemsep=2pt}
\numberwithin{equation}{section}
\theoremstyle{plain}
\newtheorem{twierdzenie}{Twierdzenie}[section]
\newtheorem{stwierdzenie}[twierdzenie]{Stwierdzenie}
\newtheorem{lemat}[twierdzenie]{Lemat}
\newtheorem{wniosek}[twierdzenie]{Wniosek}
\theoremstyle{definition}
\newtheorem{definicja}[twierdzenie]{Definicja}
\newtheorem{przyklad}[twierdzenie]{Przykład}
\theoremstyle{remark}
\newtheorem{uwaga}[twierdzenie]{Uwaga}
\newcommand{\R}{\mathbb R}
\newcommand{\C}{\mathbb C}
\newcommand{\N}{\mathbb N}
\newcommand{\K}{\mathbb K}
\newcommand{\cl}[1]{\overline{#1}}
\newcommand{\TVS}{\mathrm{TVS}}
\DeclareMathOperator{\lin}{lin}
\DeclareMathOperator{\supp}{supp}
\DeclareMathOperator{\id}{id}
\title{Geometria różniczkowa w przestrzeniach funkcyjnych}
\author{}
\date{Wersja robocza migracji -- 28 listopada 2026}

% BEGIN GLOBAL CODE LISTINGS (generated from latex/pl/code-listings.tex)
\usepackage{xcolor}
\usepackage{listings}
\definecolor{codebg}{RGB}{248,249,251}
\definecolor{codeborder}{RGB}{180,191,203}
\definecolor{codecomment}{RGB}{61,113,79}
\definecolor{codekeyword}{RGB}{41,78,132}
\lstdefinelanguage{Julia}{
  morekeywords={abstract,baremodule,begin,break,catch,const,continue,do,else,elseif,
    end,export,false,finally,for,function,global,if,import,in,isa,let,local,macro,
    module,mutable,primitive,quote,return,struct,true,try,type,using,where,while,
    missing,nothing,println,print},
  sensitive=true,morecomment=[l]{\#},morecomment=[n]{\#=}{=\#},
  morestring=[b]",
}
\lstdefinestyle{skrypt}{
  backgroundcolor=\color{codebg},frame=single,rulecolor=\color{codeborder},
  framesep=5pt,frameround=tttt,framerule=.5pt,
  basicstyle=\ttfamily\footnotesize,keywordstyle=\color{codekeyword}\bfseries,
  commentstyle=\color{codecomment}\itshape,stringstyle=\color{purple!65!black},
  numbers=left,numberstyle=\tiny\color{gray!75!black},numbersep=8pt,
  stepnumber=1,breaklines=true,breakatwhitespace=false,
  columns=fullflexible,keepspaces=true,showstringspaces=false,tabsize=4,
  linewidth=.90\linewidth,captionpos=b,aboveskip=11pt,belowskip=12pt,
  xleftmargin=12pt,xrightmargin=2pt,
}
\lstdefinestyle{pythonpde}{style=skrypt,language=Python}
\lstdefinestyle{python}{style=skrypt,language=Python}
\lstdefinestyle{julia}{style=skrypt,language=Julia}
\lstset{style=skrypt}
\renewcommand{\lstlistingname}{Kod}
\renewcommand{\lstlistlistingname}{Spis kodów}
% END GLOBAL CODE LISTINGS

\externaldocument[AF-ext-]{functional_analysis}
\makeatletter
\newcommand{\importFAreference}[1]{%
  \@ifundefined{r@AF-ext-#1}{%
    \PackageWarningNoLine{function-space-geometry}{Missing AF reference #1; compile functional_analysis first}%
  }{\expandafter\let\csname r@#1\expandafter\endcsname\csname r@AF-ext-#1\endcsname}}
\makeatother
\importFAreference{sec:integral-vector15b}
\importFAreference{cor:weak-convex-closure}
\importFAreference{thm:orth-decomposition13}
\importFAreference{thm:krein-milman}
\importFAreference{cor:bauer9}
\importFAreference{cor:alaoglu-norm}
\providecommand{\atlasref}[2]{#2}
\begin{document}
\maketitle
\tableofcontents
\clearpage
% Wersja źródłowa: audyt II, jednostka 14

\section{Pochodne Gâteaux i Frécheta w przestrzeniach Banacha i Hilberta}
\label{sec:gateaux-frechet14}

\providecommand{\atlasref}[2]{#2}


\subsection{Przygotowanie: kierunki, przyrosty i liniowe przybliżenia}

\noindent\emph{Powiązania tematyczne w Bibliotece matematycznej:} \atlasref{AM1-20::ma1:thm:linear-approximation}{AM1-20: liniowe przybliżenie funkcji jednej zmiennej}.


W analizie funkcji jednej zmiennej pochodna opisuje najlepsze liniowe przybliżenie funkcji. Dla odwzorowania między przestrzeniami unormowanymi możliwe są dwie naturalne interpretacje: przybliżenie \emph{osobno w każdym ustalonym kierunku} oraz przybliżenie \emph{jednostajne względem wszystkich małych przyrostów}. Pierwsza prowadzi do pochodnej Gâteaux, druga do pochodnej Frécheta. Różnica jest istotna już w $\R^2$.

Niech $E,F$ będą rzeczywistymi przestrzeniami Banacha, $U\subseteq E$ zbiorem otwartym i $f:U\to F$. Przez $\mathcal L(E,F)$ oznaczamy przestrzeń ograniczonych operatorów liniowych, wyposażoną w normę operatorową. Otwartość $U$ zapewnia, że dla $x\in U$ i $h\in E$ punkty $x+th$ należą do $U$ dla wszystkich dostatecznie małych $|t|$.

\begin{definicja}[Pochodna kierunkowa]
\label{def:directional14}
Niech $E,F$ będą rzeczywistymi przestrzeniami unormowanymi,
$U\subseteq E$ zbiorem otwartym, $f:U\to F$, $x\in U$ i $h\in E$.
Mówimy, że $f$ ma \emph{dwustronną pochodną kierunkową} w punkcie $x$ w kierunku $h$, jeśli w $F$ istnieje granica
$\partial_h f(x)=\lim_{t\to0,\,t\in\R}(f(x+th)-f(x))/t$.
Nie wymagamy, aby zależność $h\mapsto\partial_h f(x)$ była liniowa ani ciągła. Istnienie wszystkich pochodnych kierunkowych samo w sobie nie oznacza różniczkowalności Gâteaux.
\end{definicja}

\begin{definicja}[Różniczkowalność Gâteaux]
\label{def:gateaux14}
Niech $E,F$ będą rzeczywistymi przestrzeniami unormowanymi,
$U\subseteq E$ zbiorem otwartym, $f:U\to F$ i $x\in U$.
Odwzorowanie $f$ jest \emph{różniczkowalne w sensie Gâteaux} w $x$, jeśli istnieje $A\in\mathcal L(E,F)$ takie, że dla każdego $h\in E$ zachodzi $\partial_h f(x)=Ah$. Operator $A$, jednoznacznie wyznaczony przez te granice, oznaczamy $D_Gf(x)$. Przyjęliśmy definicję wymagającą ciągłości pochodnej jako operatora liniowego; w niektórych źródłach spotyka się słabszą konwencję, dopuszczającą nieograniczone odwzorowanie liniowe.
\end{definicja}

\begin{definicja}[Różniczkowalność Frécheta]
\label{def:frechet14}
Niech $E,F$ będą rzeczywistymi przestrzeniami unormowanymi,
$U\subseteq E$ zbiorem otwartym, $f:U\to F$ i $x\in U$.
Odwzorowanie $f$ jest \emph{różniczkowalne w sensie Frécheta} w $x$, jeśli istnieje $A\in\mathcal L(E,F)$ oraz funkcja $r$ określona dla małych $h\in E$ taka, że
\[
f(x+h)=f(x)+Ah+r(h),\qquad
\lim_{\|h\|_E\to0,\,h\ne0}\frac{\|r(h)\|_F}{\|h\|_E}=0.
\]
Operator $A$ oznaczamy $Df(x)$ albo $D_Ff(x)$; piszemy również $r(h)=o(\|h\|_E)$. Symbol $Df(x)h$ oznacza wartość operatora $Df(x)$ na wektorze $h$, nie jego iloczyn skalarny z $h$.
\end{definicja}

\begin{uwaga}[Założenia o zupełności]
\label{rem:completeness14}
Do samych definicji wystarczą przestrzenie unormowane i otwarta dziedzina. Zupełność potrzebna będzie między innymi do twierdzenia o całkowaniu funkcji o wartościach w przestrzeni Banacha oraz w późniejszych twierdzeniach o funkcji odwrotnej. Stąd w niniejszym rozdziale pracujemy przede wszystkim w przestrzeniach Banacha.
\end{uwaga}

\subsection{Implikacje między pochodnymi}

\begin{twierdzenie}[Fréchet implikuje Gâteaux]
\label{thm:frechet-implies-gateaux14}
Niech $E,F$ będą rzeczywistymi przestrzeniami unormowanymi, $U\subseteq E$ zbiorem otwartym, a $f:U\to F$ odwzorowaniem różniczkowalnym w sensie Frécheta w punkcie $x\in U$. Wówczas $f$ jest ciągłe w $x$, jest różniczkowalne w sensie Gâteaux w $x$ oraz $D_Gf(x)=Df(x)$. Pochodna Frécheta, jeśli istnieje, jest jednoznaczna.
\end{twierdzenie}
\begin{proof}
Niech $A=Df(x)$ i $f(x+h)-f(x)=Ah+r(h)$, gdzie $\|r(h)\|=o(\|h\|)$. Ponieważ $A$ jest ograniczony, $\|Ah+r(h)\|\le\|A\|\|h\|+o(\|h\|)\to0$, co daje ciągłość. Dla ustalonego $h\ne0$ i małego rzeczywistego $t$ mamy
$\|(f(x+th)-f(x))/t-Ah\|=\|r(th)\|/|t|=\|h\|\,\|r(th)\|/\|th\|\to0$.
Dla $h=0$ pochodna kierunkowa wynosi zero. Jeżeli $A$ i $B$ spełniają definicję Frécheta, powyższy argument daje $Ah=Bh$ dla każdego $h$, a więc $A=B$.
\end{proof}

\begin{stwierdzenie}[Kryterium jednostajnego przybliżenia kierunkowego]
\label{prop:uniform-directions14}
Niech $E,F$ będą rzeczywistymi przestrzeniami unormowanymi, $U\subseteq E$ otwarty, $x\in U$ i $A\in\mathcal L(E,F)$. Odwzorowanie $f:U\to F$ ma pochodną Frécheta $Df(x)=A$ wtedy i tylko wtedy, gdy
\[
\lim_{t\to0,\,t\in\R}\ \sup_{\|h\|_E=1}
\left\|\frac{f(x+th)-f(x)}{t}-Ah\right\|_F=0.
\]
Zakładamy $E\ne\{0\}$. Przez $B_E$ rozumiemy domkniętą kulę jednostkową w $E$; supremum rozpatrujemy wyłącznie dla $t\ne0$ o module na tyle małym, że $x+tB_E\subseteq U$.
\end{stwierdzenie}
\begin{proof}
Dla reszty $r(k)=f(x+k)-f(x)-Ak$ mamy
$\sup_{\|h\|=1}\|r(th)\|/|t|=\sup_{\|k\|=|t|}\|r(k)\|/\|k\|$.
Warunek Frécheta zapewnia, że te suprema dążą do zera. Odwrotnie, dla $k\ne0$ połóżmy $t=\|k\|$ i $h=k/\|k\|$; uzyskana jednostajna granica daje $\|r(k)\|/\|k\|\to0$.
\end{proof}

\begin{uwaga}
\label{rem:quantifiers14}
W definicji Gâteaux kwantyfikatory mają kolejność: \emph{dla każdego ustalonego $h$ granica przy $t\to0$}. W definicji Frécheta wymagamy wspólnej kontroli błędu dla wszystkich kierunków. Zamiana kolejności kwantyfikatorów nie jest uprawniona, nawet w skończonym wymiarze.
\end{uwaga}

\subsection{Kontrprzykłady: czego pochodna Gâteaux nie gwarantuje}

\begin{przyklad}[Wszystkie pochodne kierunkowe i brak ciągłości]
\label{ex:directional-discontinuous14}
Dla $(x,y)\ne(0,0)$ połóżmy
$g(x,y)=x^3y/(x^6+y^2)$, zaś $g(0,0)=0$. Dla każdego $(a,b)\in\R^2$ granica $\partial_{(a,b)}g(0,0)$ istnieje i jest równa $0$: jeśli $b\ne0$, to $g(ta,tb)/t=a^3b\,t/(a^6t^4+b^2)\to0$; jeśli $b=0$, iloraz jest zerowy. Zatem $D_Gg(0,0)=0\in\mathcal L(\R^2,\R)$. Jednak na krzywej $y=x^3$, $x\ne0$, mamy $g(x,x^3)=1/2$. Funkcja nie jest ciągła w $(0,0)$, więc nie jest tam różniczkowalna w sensie Frécheta.
\end{przyklad}

\begin{przyklad}[Ciągłość i pochodna Gâteaux nie implikują pochodnej Frécheta]
\label{ex:continuous-gateaux-not-frechet14}
Niech $f(0,0)=0$, zaś dla $(x,y)\ne(0,0)$ niech
\[
f(x,y)=\frac{x^3y}{x^4+y^2}.
\]
Z nierówności $2|x^2y|\le x^4+y^2$ wynika $|f(x,y)|\le|x|/2$; zatem $f$ jest ciągła w $(0,0)$. Dla ustalonego $(a,b)$, gdy $b\ne0$, otrzymujemy
$f(ta,tb)/t=a^3b\,t/(a^4t^2+b^2)\to0$, a dla $b=0$ iloraz jest zerowy. Stąd $D_Gf(0,0)=0$. Natomiast $f(x,x^2)=x/2$ dla $x\ne0$ oraz
$|f(x,x^2)|/\sqrt{x^2+x^4}=1/(2\sqrt{1+x^2})\to1/2$ przy $x\to0^+$. Reszta względem zerowego operatora nie jest $o(\|(x,y)\|_2)$: pochodna Frécheta nie istnieje.
\end{przyklad}

\begin{przyklad}[Sama pochodna kierunkowa nie musi być liniowa]
\label{ex:nonlinear-directional14}
Dla $f(x,y)=|x|+|y|$ w zerze istnieje tylko jednostronna pochodna kierunkowa, równa $|a|+|b|$ w kierunku $(a,b)$; dwustronna pochodna kierunkowa przy $h\ne0$ nie istnieje. Dla przykładu z \emph{dwustronnymi} pochodnymi kierunkowymi, lecz nieliniową zależnością od kierunku, rozważmy $q(0,0)=0$ oraz $q(x,y)=x^3/(x^2+y^2)$ dla $(x,y)\ne0$. Wówczas $\partial_{(a,b)}q(0,0)=a^3/(a^2+b^2)$ dla $(a,b)\ne(0,0)$, a w kierunku zerowym pochodna wynosi $0$. Odwzorowanie to jest jednorodne względem mnożenia kierunku przez rzeczywiste skalary, lecz nie jest addytywne: $\partial_{(1,0)}q+\partial_{(0,1)}q=1$, podczas gdy $\partial_{(1,1)}q=1/2$. Dlatego $q$ nie ma pochodnej Gâteaux według Definicji~\ref{def:gateaux14}.
\end{przyklad}

\begin{stwierdzenie}[Brak reguły łańcuchowej dla samej pochodnej Gâteaux]
\label{prop:chain-fails14}
Istnieją ciągłe odwzorowanie $f:\R^2\to\R$, różniczkowalne w sensie Gâteaux w zerze, oraz gładka krzywa $\gamma:\R\to\R^2$ z $\gamma(0)=0$ takie, że $\frac{\mathrm d}{\mathrm dt}(f\circ\gamma)(0)\ne D_Gf(0)\gamma'(0)$.
\end{stwierdzenie}
\begin{proof}
Niech $f$ będzie funkcją z Przykładu~\ref{ex:continuous-gateaux-not-frechet14}, a $\gamma(t)=(t,t^2)$. Wówczas $f(\gamma(t))=t/2$ dla $t\ne0$, a także dla $t=0$ z ciągłości. Zatem $(f\circ\gamma)'(0)=1/2$, podczas gdy $D_Gf(0)=0$. Niezgodność wynika z niestabilności kierunkowego przybliżenia przy zmiennym kierunku przyrostu.
\end{proof}

\begin{center}
\begin{tikzpicture}[x=2.15cm,y=1.65cm,>=Stealth]
  \draw[->,thin] (-0.17,0)--(1.7,0) node[below] {$x$};
  \draw[->,thin] (0,-0.42)--(0,1.48) node[left] {$y$};
  \draw[gray,densely dashed] (0,0)--(1.32,0.62);
  \draw[gray,densely dashed] (0,0)--(0.79,1.22);
  \draw[gray,densely dashed] (0,0)--(1.42,-0.29);
  \draw[very thick,blue!55!black,domain=0:1.15,samples=70,smooth]
       plot (\x,{\x*\x});
  \fill (0,0) circle (1.05pt);
  \node[anchor=west,blue!55!black] at (1.17,1.22) {\footnotesize $y=x^2$};
  \node[anchor=west] at (0.83,-0.37) {\scriptsize kierunki stałe};
\end{tikzpicture}

{\small Proste przez zero odpowiadają ustalonym kierunkom w definicji
Gâteaux. Krzywa $y=x^2$ prowadzi do zera po zmieniającym się
kierunku; na niej ujawnia się brak różniczkowalności Frécheta.}
\end{center}

\subsection{Pochodna Gâteaux normy w \texorpdfstring{$\ell^1$}{ell1} bez pochodnej Frécheta}

\begin{twierdzenie}[Kontrprzykład w nieskończonym wymiarze]
\label{thm:l1-gateaux14}
Niech $E=\ell^1(\N;\R)$ i $N:E\to\R$, $N(x)=\|x\|_1$. Jeśli $x=(x_n)\in E$ spełnia $x_n\ne0$ dla wszystkich $n\ge1$, to $N$ jest różniczkowalna w sensie Gâteaux w $x$ i dla każdego $h=(h_n)\in E$
$D_GN(x)h=\sum_{n=1}^\infty \operatorname{sgn}(x_n)h_n$.
Funkcja $N$ nie jest jednak różniczkowalna w sensie Frécheta w żadnym takim punkcie $x$.
\end{twierdzenie}
\begin{proof}
Dla ustalonego $h\in\ell^1$ i $t\ne0$ oznaczmy $a_n(t)=(|x_n+th_n|-|x_n|)/t$. Ponieważ $x_n\ne0$, mamy $a_n(t)\to\operatorname{sgn}(x_n)h_n$ dla każdego $n$, przy czym $|a_n(t)|\le|h_n|$. Zbieżność zmajoryzowana dla szeregu (albo argument z obcięciem sumy i oszacowaniem ogona przez $\sum_{n>N}|h_n|$) daje
$\lim_{t\to0}(N(x+th)-N(x))/t=\sum_n \operatorname{sgn}(x_n)h_n$.
Szereg jest bezwzględnie zbieżny, a liniowość wynika z liniowości sum częściowych i przejścia do granicy. Moduł sumy nie przekracza $\|h\|_1$, więc funkcjonał jest ciągły i ma normę co najwyżej $1$. Dla $h=\operatorname{sgn}(x_1)e_1$ przyjmuje wartość $1$ na wektorze o normie $1$, zatem jego norma wynosi dokładnie $1$.

Jeżeli istniałaby pochodna Frécheta, musiałaby być równa $D_GN(x)$. Dla $h^{(n)}=-2x_ne_n$ mamy $\|h^{(n)}\|_1=2|x_n|\to0$, ponieważ $x\in\ell^1$, a zarazem $N(x+h^{(n)})=N(x)$ i $D_GN(x)h^{(n)}=-2|x_n|$. Wobec tego
\[
\frac{|N(x+h^{(n)})-N(x)-D_GN(x)h^{(n)}|}{\|h^{(n)}\|_1}=1
\]
dla każdego $n$. Jest to sprzeczne z warunkiem Frécheta.
\end{proof}

\begin{uwaga}
Jeśli pewna współrzędna $x_n$ znika, norma $\ell^1$ nie ma w tym punkcie nawet dwustronnej pochodnej kierunkowej w kierunku $e_n$. Powyższy kontrprzykład pokazuje więc ścisłą różnicę między Gâteaux a Fréchetem również w punktach, w których wszystkie pochodne kierunkowe tworzą ciągły funkcjonał liniowy.
\end{uwaga}

\subsection{Reguły różniczkowania w przestrzeniach Banacha}

\begin{stwierdzenie}[Suma i operator liniowy]
\label{prop:sum-linear14}
Niech $E,F,G$ będą rzeczywistymi przestrzeniami Banacha, $U\subseteq E$ otwarty i $x\in U$. Jeżeli $f,g:U\to F$ są różniczkowalne w sensie Frécheta w $x$, a $\alpha,\beta\in\R$, to $D(\alpha f+\beta g)(x)=\alpha Df(x)+\beta Dg(x)$. Jeśli $T\in\mathcal L(F,G)$, to $D(T\circ f)(x)=T\circ Df(x)$. Ponadto dla $L\in\mathcal L(E,F)$ i $b\in F$ funkcja $f(y)=Ly+b$ ma $Df(y)=L$ w każdym $y\in E$.
\end{stwierdzenie}
\begin{proof}
Wstawiamy rozwinięcia $f(x+h)=f(x)+Df(x)h+r_f(h)$ i $g(x+h)=g(x)+Dg(x)h+r_g(h)$, gdzie normy obu reszt są $o(\|h\|)$. Dla kombinacji liniowej otrzymujemy ograniczony operator $\alpha Df(x)+\beta Dg(x)$ oraz resztę $\alpha r_f(h)+\beta r_g(h)$, której norma podzielona przez $\|h\|$ jest nie większa niż $|\alpha|\|r_f(h)\|/\|h\|+|\beta|\|r_g(h)\|/\|h\|\to0$. Dla złożenia z $T$ operator jest równy $TDf(x)$, a reszta $Tr_f(h)$ spełnia $\|Tr_f(h)\|\le\|T\|\|r_f(h)\|=o(\|h\|)$. W przypadku funkcji afinicznej $f(y+h)-f(y)=Lh$, więc reszta wynosi zero.
\end{proof}

\noindent\emph{Powiązanie tematyczne w Bibliotece matematycznej:} \atlasref{DG-12::def:pochodna-wzdluz-krzywej}{DG-12: zastosowanie reguły łańcuchowej do pochodnej wzdłuż krzywej}.

\begin{twierdzenie}[Reguła łańcuchowa Frécheta]
\label{thm:chain14}
Niech $E,F,G$ będą rzeczywistymi przestrzeniami Banacha, $U\subseteq E$, $V\subseteq F$ otwarte, $f:U\to F$ oraz $g:V\to G$ i $f(U)\subseteq V$. Jeśli $f$ jest różniczkowalne w sensie Frécheta w $x\in U$, a $g$ w $f(x)$, to $g\circ f$ jest różniczkowalne w sensie Frécheta w $x$, a
$D(g\circ f)(x)=Dg(f(x))\circ Df(x)$.
\end{twierdzenie}
\begin{proof}
Oznaczmy $A=Df(x)$, $B=Dg(f(x))$. Wówczas $f(x+h)=f(x)+Ah+r(h)$, $\|r(h)\|=o(\|h\|)$, i dla $k\to0$ mamy $g(f(x)+k)=g(f(x))+Bk+s(k)$, gdzie $\|s(k)\|=o(\|k\|)$. Ponieważ $f$ jest ciągłe w $x$, dla małego $h$ można przyjąć $k=Ah+r(h)$, a $\|k\|\le(\|A\|+1)\|h\|$. Otrzymujemy
$g(f(x+h))-g(f(x))-BAh=Br(h)+s(Ah+r(h))$.
Pierwszy składnik jest $o(\|h\|)$ przez ograniczoność $B$. Drugi także: jeśli $k\ne0$, to
$\|s(k)\|/\|h\|=(\|s(k)\|/\|k\|)(\|k\|/\|h\|)\to0$;
jeśli $k=0$, to $s(k)=0$. Dowód nie wymaga ciągłości pochodnych w otoczeniu rozważanych punktów.
\end{proof}

\begin{twierdzenie}[Reguła dla odwzorowania dwuliniowego]
\label{thm:bilinear14}
Niech $E,E_1,E_2,G$ będą rzeczywistymi przestrzeniami Banacha i $B:E_1\times E_2\to G$ ciągłym odwzorowaniem dwuliniowym, tj. $\|B(u,v)\|\le C\|u\|\|v\|$. Wówczas $B$ jest różniczkowalne w sensie Frécheta na przestrzeni produktowej $E_1\times E_2$ i
$DB(u,v)(h,k)=B(h,v)+B(u,k)$.
Jeżeli $f:U\to E_1$, $g:U\to E_2$ są różniczkowalne w $x\in U\subseteq E$ otwartym, to
$D(B(f,g))(x)h=B(Df(x)h,g(x))+B(f(x),Dg(x)h)$.
\end{twierdzenie}
\begin{proof}
Dokładna tożsamość
$B(u+h,v+k)-B(u,v)-B(h,v)-B(u,k)=B(h,k)$
daje $\|B(h,k)\|\le C\|h\|\|k\|\le C\max(\|h\|,\|k\|)^2$. Reszta jest więc $o(\max(\|h\|,\|k\|))$, a operator liniowy wskazany w tezie jest ograniczony. Dla odwzorowania $z\mapsto(f(z),g(z))$ pochodna jest równa $h\mapsto(Df(x)h,Dg(x)h)$: w normie maksimum na produkcie iloraz normy reszty przez $\|h\|$ jest maksimum dwóch ilorazów reszt, które dążą do zera. Operator ten jest ograniczony przez $\max(\|Df(x)\|,\|Dg(x)\|)$. Reguła łańcuchowa zastosowana do tego odwzorowania i do $B$ daje ostatni wzór.
\end{proof}

\begin{przyklad}[Operatory i funkcje kwadratowe]
\label{ex:quadratic14}
Niech $E$ będzie rzeczywistą przestrzenią Banacha, $F$ rzeczywistą przestrzenią Banacha i $B:E\times E\to F$ ciągłym odwzorowaniem dwuliniowym. Dla $Q:E\to F$, $Q(x)=B(x,x)$, oraz każdego $x,h\in E$ zachodzi
$DQ(x)h=B(h,x)+B(x,h)$.
W szczególności, jeśli $H$ jest rzeczywistą przestrzenią Hilberta i $\Phi(x)=\frac12\|x\|^2$, to
$D\Phi(x)h=\langle h,x\rangle$, a reszta ma postać $\frac12\|h\|^2$.
\end{przyklad}

\subsection{Różniczkowalność klasy \texorpdfstring{$C^1$}{C1} i twierdzenie o przyrostach}

\begin{definicja}[Klasa $C^1$]
\label{def:c1-14}
Niech $E,F$ będą rzeczywistymi przestrzeniami Banacha, $U\subseteq E$ otwarty i $f:U\to F$. Piszemy $f\in C^1(U,F)$, jeśli $f$ jest różniczkowalne w sensie Frécheta w każdym punkcie $U$ oraz odwzorowanie $Df:U\to\mathcal L(E,F)$ jest ciągłe w normie operatorowej.
\end{definicja}

\begin{lemat}[Podstawowe twierdzenie rachunku dla krzywych Banacha]
\label{lem:ftc-banach14}
Niech $F$ będzie rzeczywistą przestrzenią Banacha, $a<b$ i $\gamma:[a,b]\to F$ funkcją klasy $C^1$ (pochodna w końcach jednostronna i ciągle przedłużona). Wówczas $\gamma'$ jest całkowalna w sensie Riemanna z wartościami w $F$ oraz
$\gamma(b)-\gamma(a)=\int_a^b\gamma'(t)\,\mathrm dt$.
\end{lemat}
\begin{proof}
Ciągłość $\gamma'$ na zwartym przedziale daje jej jednostajną ciągłość. Dla dwóch znakowanych podziałów o średnicach mniejszych od $\delta$ wybieramy wspólne rozdrobnienie. Różnica sumy na każdym z pierwotnych podziałów i sumy na rozdrobnieniu ma normę co najwyżej $(b-a)\omega(\delta)$, gdzie $\omega(\delta)=\sup_{|s-t|\le\delta}\|\gamma'(s)-\gamma'(t)\|$ dla $s,t\in[a,b]$. Istotnie, znak podprzedziału rozdrobnienia i znak zawierającego go przedziału leżą w tym samym przedziale długości mniejszej niż $\delta$, a długości podprzedziałów sumują się do $b-a$. Stąd różnica dwóch pierwotnych sum ma normę co najwyżej $2(b-a)\omega(\delta)\to0$. Wybierając ciąg podziałów o średnicach dążących do zera, otrzymujemy ciąg Cauchy'ego sum, który ma granicę dzięki zupełności $F$. To samo oszacowanie pokazuje, że wszystkie sumy dostatecznie drobnych podziałów dążą do tej granicy, czyli całka Riemanna istnieje. Dla podziału $a=t_0<\cdots<t_m=b$ wybierzmy dowolne $\xi_j\in[t_{j-1},t_j]$. Z definicji pochodnej i jednostajnej ciągłości $\gamma'$ wynika, że przy dostatecznie małej średnicy podziału
$\|\gamma(t_j)-\gamma(t_{j-1})-(t_j-t_{j-1})\gamma'(\xi_j)\|\le\varepsilon(t_j-t_{j-1})$.
Ściślej, dla dowolnego rzeczywistego funkcjonału $\ell\in F'$ o normie co najwyżej $1$ stosujemy skalarne twierdzenie o wartości średniej do $t\mapsto\ell(\gamma(t)-t\gamma'(\xi_j))$ na $[t_{j-1},t_j]$. Otrzymujemy górne oszacowanie modułu przyrostu przez $(t_j-t_{j-1})\sup_{t\in[t_{j-1},t_j]}\|\gamma'(t)-\gamma'(\xi_j)\|$. Z twierdzenia Hahna--Banacha, po supremum po takich $\ell$, dostajemy to samo oszacowanie normy wektora. Jednostajna ciągłość $\gamma'$ zapewnia żądaną nierówność dla dostatecznie drobnych podziałów. Sumowanie oraz przejście z średnicą podziału do zera kończy dowód.
\end{proof}

\begin{twierdzenie}[Wzór całkowy i oszacowanie przyrostu]
\label{thm:mean-value14}
Niech $E,F$ będą rzeczywistymi przestrzeniami Banacha, $U\subseteq E$ zbiorem otwartym, $f\in C^1(U,F)$, a $x,y\in U$ takimi punktami, że $[x,y]=\{x+t(y-x):t\in[0,1]\}\subseteq U$. Wówczas
\[
f(y)-f(x)=\int_0^1 Df(x+t(y-x))(y-x)\,\mathrm dt,
\]
gdzie całka ma wartości w $F$. W szczególności, jeśli $\sup_{z\in[x,y]}\|Df(z)\|\le M$, to $\|f(y)-f(x)\|\le M\|y-x\|$.
\end{twierdzenie}
\begin{proof}
Definiujemy $\gamma(t)=f(x+t(y-x))$. Reguła łańcuchowa daje $\gamma'(t)=Df(x+t(y-x))(y-x)$; ciągłość $Df$ zapewnia ciągłość $\gamma'$. Lemat~\ref{lem:ftc-banach14} daje wzór. Oszacowanie normy całki przez całkę normy wynika z oszacowania sum Riemanna i przejścia do granicy.
\end{proof}

\begin{twierdzenie}[Ciągła pochodna Gâteaux daje pochodną Frécheta]
\label{thm:continuous-gateaux14}
Niech $E,F$ będą rzeczywistymi przestrzeniami Banacha, $U\subseteq E$ otwarty, $f:U\to F$. Załóżmy, że $f$ jest różniczkowalne w sensie Gâteaux w każdym punkcie $U$ i odwzorowanie $D_Gf:U\to\mathcal L(E,F)$ jest ciągłe w normie operatorowej. Wtedy $f\in C^1(U,F)$ oraz $Df(x)=D_Gf(x)$ dla wszystkich $x\in U$.
\end{twierdzenie}
\begin{proof}
Ustalmy $x\in U$ i weźmy małe $h$ takie, że $[x,x+h]\subseteq U$. Funkcja $\gamma(t)=f(x+th)$ jest różniczkowalna dla $t\in(0,1)$ i $\gamma'(t)=D_Gf(x+th)h$, co wynika z definicji kierunkowej. Jej pochodna jest ciągła dzięki założeniu, a $\gamma$ jest ciągła również w końcach przedziału, ponieważ istnienie pochodnej kierunkowej w tych punktach zapewnia ciągłość wzdłuż tej prostej. Z Lematu~\ref{lem:ftc-banach14} mamy
$f(x+h)-f(x)=\int_0^1 D_Gf(x+th)h\,\mathrm dt$.
Po odjęciu $D_Gf(x)h$ i oszacowaniu normy dostajemy
\[
\frac{\|f(x+h)-f(x)-D_Gf(x)h\|}{\|h\|}
\le\sup_{0\le t\le1}\|D_Gf(x+th)-D_Gf(x)\|\longrightarrow0.
\]
To dowodzi różniczkowalności Frécheta; ciągłość pochodnej otrzymujemy z założenia.
\end{proof}

\begin{przyklad}[Pochodna istniejąca, ale nieciągła]
\label{ex:derivative-noncontinuous14}
Na $\R$ niech $f(0)=0$ oraz $f(x)=x^2\sin(1/x^2)$ dla $x\ne0$. Ponieważ $|f(h)/h|\le|h|$, mamy $f'(0)=0$. Dla $x\ne0$ reguła iloczynu i reguła łańcuchowa dają $f'(x)=2x\sin(1/x^2)-2x^{-1}\cos(1/x^2)$. Dla $x_n=(2\pi n)^{-1/2}$ zachodzi $x_n\to0$ i $f'(x_n)=-2/x_n$, więc pochodna nie jest ciągła w zerze. Zatem $f$ jest różniczkowalna Frécheta w każdym punkcie, lecz nie należy do $C^1(\R,\R)$. Nie należy utożsamiać istnienia pochodnej z jej ciągłością.
\end{przyklad}

\subsection{Norma i gradient w rzeczywistych i zespolonych przestrzeniach Hilberta}

\begin{twierdzenie}[Różniczkowanie kwadratu normy i normy]
\label{thm:hilbert-norm14}
Niech $H$ będzie przestrzenią Hilberta nad $\K\in\{\R,\C\}$, traktowaną jako \emph{rzeczywista} przestrzeń Banacha, i niech $\Phi(x)=\frac12\|x\|^2$. W każdym $x\in H$ istnieje rzeczywista pochodna Frécheta $D_{\R}\Phi(x)h=\operatorname{Re}\langle h,x\rangle$. Dla $x\ne0$ norma $N(x)=\|x\|$ ma rzeczywistą pochodną Frécheta
$D_{\R}N(x)h=\operatorname{Re}\langle h,x\rangle/\|x\|$.
Jeżeli $H\ne\{0\}$, norma nie jest różniczkowalna Gâteaux w $0$ jako funkcja na rzeczywistej przestrzeni Banacha.
\end{twierdzenie}
\begin{proof}
Z rozwinięcia iloczynu skalarnego wynika dokładna tożsamość
$\Phi(x+h)-\Phi(x)-\operatorname{Re}\langle h,x\rangle=\frac12\|h\|^2$, a funkcjonał $h\mapsto\operatorname{Re}\langle h,x\rangle$ jest rzeczywisto-liniowy i ograniczony. To dowodzi pierwszej tezy. Dla $x\ne0$ stosujemy rzeczywistą regułę łańcuchową do $N(x)=\sqrt{2\Phi(x)}$, ponieważ pierwiastek jest gładki na $(0,\infty)$. W zerze iloraz kierunkowy wynosi $\|th\|/t=(|t|/t)\|h\|$; dla $h\ne0$ nie istnieje granica dwustronna.
\end{proof}

\begin{definicja}[Gradient na rzeczywistej przestrzeni Hilberta]
\label{def:gradient14}
Niech $H$ będzie rzeczywistą przestrzenią Hilberta, $U\subseteq H$ otwarty i $f:U\to\R$ różniczkowalna w sensie Frécheta w $x\in U$. \emph{Gradientem} $f$ w $x$ nazywamy jedyny wektor $\nabla f(x)\in H$ spełniający $Df(x)h=\langle h,\nabla f(x)\rangle$ dla wszystkich $h\in H$. Istnienie i jednoznaczność daje twierdzenie Riesza. Dla zespolonej przestrzeni Hilberta i funkcjonału rzeczywistego używamy struktury rzeczywistej i iloczynu skalarnego $(h,k)_{\R}=\operatorname{Re}\langle h,k\rangle$.
\end{definicja}

\begin{przyklad}[Gradient energii kwadratowej]
\label{ex:quadratic-energy14}
Niech $H$ będzie rzeczywistą lub zespoloną przestrzenią Hilberta, $A\in\mathcal L(H)$ operatorem samosprzężonym, $b\in H$, a $J(x)=\frac12\langle Ax,x\rangle-\operatorname{Re}\langle x,b\rangle$. Z samosprzężoności wynika, że $J$ ma wartości rzeczywiste. Dla $h\in H$ otrzymujemy
$D_{\R}J(x)h=\operatorname{Re}\langle h,Ax-b\rangle$, więc $\nabla J(x)=Ax-b$ w rzeczywistej strukturze Hilberta. Jeżeli $A$ spełnia oszacowanie $\langle Ax,x\rangle\ge\alpha\|x\|^2$ dla pewnego $\alpha>0$, to z nierówności Cauchy'ego--Schwarza wynika $\|Ax\|\ge\alpha\|x\|$. Operator jest iniektywny, a jego obraz jest domknięty: jeśli $Ax_n$ zbiega, to $x_n$ jest ciągiem Cauchy'ego i z ciągłości $A$ granica $Ax_n$ należy do obrazu. Ponadto wektor $y$ prostopadły do obrazu spełnia $\langle x,Ay\rangle=\langle Ax,y\rangle=0$ dla każdego $x$, więc $Ay=0$ i $y=0$. Obraz jest zatem gęsty i domknięty, czyli równy $H$. Istnieje dokładnie jeden $x_*$ z $Ax_*=b$. Rozwinięcie formy kwadratowej i samosprzężoność dają dla każdego $h\in H$
$J(x_*+h)-J(x_*)=\frac12\langle Ah,h\rangle\ge\frac\alpha2\|h\|^2$.
Zatem $x_*$ jest jedynym minimizatorem $J$.
\end{przyklad}

\begin{przyklad}[Funkcjonał całkowy i pochodna na $L^2$]
\label{ex:integral-functional14}
Niech $(\Omega,\mathcal A,\mu)$ będzie przestrzenią z miarą skończoną i $H=L^2(\Omega,\mu;\R)$. Dla $g\in H$ określmy $J(u)=\frac12\int_\Omega |u(t)-g(t)|^2\,\mathrm d\mu(t)$. Wtedy
$DJ(u)h=\int_\Omega (u(t)-g(t))h(t)\,\mathrm d\mu(t)$,
a $\nabla J(u)=u-g$. Wynik można otrzymać przez rozwinięcie kwadratu oraz nierówność Cauchy'ego--Schwarza; reszta ma dokładnie postać $\frac12\|h\|_2^2$.
\end{przyklad}

\subsection{Rzeczywista a zespolona różniczkowalność}

\begin{definicja}[Pochodna zespolona Frécheta i Gâteaux]
\label{def:complex-differentiability14}
Niech $E,F$ będą zespolonymi przestrzeniami Banacha, $U\subseteq E$ zbiorem otwartym, $f:U\to F$ oraz $x\in U$. Przez $\mathcal L_{\C}(E,F)$ rozumiemy przestrzeń ograniczonych operatorów zespolenie liniowych. \emph{Zespoloną pochodną Frécheta} $D_{\C}f(x)$ nazywamy operator $A\in\mathcal L_{\C}(E,F)$ spełniający $f(x+h)=f(x)+Ah+r(h)$ dla dostatecznie małych $h\in E$, gdzie $\|r(h)\|/\|h\|\to0$ przy $h\to0$, $h\ne0$. \emph{Zespolona pochodna Gâteaux} w $x$ oznacza istnienie operatora $A\in\mathcal L_{\C}(E,F)$ takiego, że dla każdego $h\in E$
$\lim_{z\to0,\,z\in\C}(f(x+zh)-f(x))/z=Ah$.
W obu definicjach liniowość wymagana jest nad $\C$, a dla pochodnej Gâteaux parametr $z$ przebiega płaszczyznę zespoloną, nie tylko oś rzeczywistą.
\end{definicja}

\begin{stwierdzenie}[Pochodna zespolona a rzeczywista]
\label{prop:complex-real14}
Niech $E,F$ będą zespolonymi przestrzeniami Banacha, $U\subseteq E$ otwarty i $f:U\to F$. Jeśli $f$ jest zespolenie różniczkowalne w sensie Frécheta w $x\in U$, to jest też rzeczywiście różniczkowalne w sensie Frécheta w $x$, a $D_{\R}f(x)$ jest operatorem rzeczywistym odpowiadającym $D_{\C}f(x)$. Odwrotnie, jeśli $D_{\R}f(x)$ istnieje, to $D_{\C}f(x)$ istnieje wtedy i tylko wtedy, gdy $D_{\R}f(x)(ih)=iD_{\R}f(x)h$ dla wszystkich $h\in E$.
\end{stwierdzenie}
\begin{proof}
Pochodna zespolona spełnia oszacowanie reszty wymagane dla pochodnej rzeczywistej, ponieważ normy nie zmieniają się po zapomnieniu struktury zespolonej. Jeśli istnieje $D_{\R}f(x)=A$, to identyczna reszta $o(\|h\|)$ daje zespoloną pochodną dokładnie wtedy, gdy rzeczywisto-liniowy operator $A$ jest zespolenie liniowy. Dla operatora rzeczywisto-liniowego $A$ warunek $A(ih)=iA(h)$ jest równoważny liniowości zespolonej.
\end{proof}

\begin{przyklad}[Sprzężenie zespolone i kwadrat modułu]
\label{ex:complex-nonholomorphic14}
Funkcja $f:\C\to\C$, $f(z)=\overline z$, jest rzeczywiście różniczkowalna Frécheta w każdym punkcie, z $D_{\R}f(z)h=\overline h$, ale nie jest zespolenie różniczkowalna w żadnym punkcie, ponieważ $h\mapsto\overline h$ nie spełnia warunku liniowości zespolonej z poprzedniego stwierdzenia (porównujemy $h=1$ i $h=i$). Podobnie $z\mapsto|z|^2$ ma rzeczywistą pochodną $h\mapsto2\operatorname{Re}(h\overline z)$, nie ma zaś zespolonej pochodnej Frécheta w punktach $z\ne0$: niezerowy operator rzeczywisto-liniowy o wartościach rzeczywistych nie może być liniowy nad $\C$. W punkcie $z=0$ zespolona pochodna Frécheta funkcji $z\mapsto|z|^2$ istnieje i wynosi zero, gdyż $|h|^2=o(|h|)$.
\end{przyklad}


\subsection{Słabe pochodne krzywych i pochodne w topologii słabej-*}

\begin{definicja}[Pochodne krzywych w różnych topologiach]
\label{def:weak-curve-derivative14}
Niech $X,Y$ będą przestrzeniami Banacha nad $\K\in\{\R,\C\}$,
$I\subseteq\R$ otwartym przedziałem, $t\in I$ oraz
$\gamma:I\to X$, $\eta:I\to Y'$ krzywymi (nie zakładamy
z góry ich ciągłości). Krzywa $\gamma$ ma w $t$ pochodną słabą $v\in X$,
jeżeli dla każdego $x'\in X'$
\[
 \lim_{h\to0}\frac{x'(\gamma(t+h))-x'(\gamma(t))}{h}=x'(v).
\]
Krzywa $\eta$ ma w $t$ pochodną słabą-* $w\in Y'$,
jeżeli dla każdego $y\in Y$
\[
 \lim_{\substack{h\to0\\h\ne0}}
 \frac{\eta(t+h)(y)-\eta(t)(y)}{h}=w(y).
\]
W obu granicach rozpatrujemy przyrosty rzeczywiste $h\ne0$
takie, że $t+h\in I$. Pochodna normowa jest granicą
ilorazów różnicowych w normie; implikuje odpowiednią
pochodną słabą lub słabą-* z ciągłości wszystkich testów.
Rozdzielanie punktów przez funkcjonały zapewnia jednoznaczność.
\end{definicja}

\begin{twierdzenie}[Istnienie pochodnej słabej-* i pochodnej w przestrzeni podwójnie sprzężonej]
\label{thm:scalar-curve-derivative14}
Niech $X,Y$ będą przestrzeniami Banacha nad $\K\in\{\R,\C\}$,
$I\subseteq\R$ otwartym przedziałem, $t\in I$,
$\gamma:I\to X$ oraz $\eta:I\to Y'$ krzywymi.
Jeżeli dla każdego $x'\in X'$ krzywa skalarna $x'\circ\gamma$ ma
pochodną w $t$, to istnieje dokładnie jeden $D\in X''$
taki, że $D(x')=(x'\circ\gamma)'(t)$. Pochodna słaba
$\gamma$ istnieje wtedy i tylko wtedy, gdy $D\in J_X(X)$,
gdzie $J_X(x)(x')=x'(x)$. W szczególności wystarcza
refleksywność $X$. Jeżeli dla każdego $y\in Y$ funkcja
$s\mapsto\eta(s)(y)$ ma pochodną w $t$, to istnieje
pochodna słaba-* $\eta$ w $t$.
\end{twierdzenie}
\begin{proof}
Wybierzmy ciąg niezerowych $h_n\to0$ z $t+h_n\in I$.
Ponieważ $X'$ jest przestrzenią Banacha, rodzina $J_X((\gamma(t+h_n)-\gamma(t))/h_n)$ jest
punktowo ograniczona na $X'$ dzięki istnieniu granic
skalarnych. Twierdzenie Banacha--Steinhausa daje
jednostajne ograniczenie jej norm. Granica punktowa
jest liniowa i ograniczona przez tę samą stałą,
a więc należy do $X''$. Jej wartości nie zależą od
wybranego ciągu, ponieważ skalarne pochodne istnieją.
Jeżeli $D=J_X(v)$, definicja daje pochodną słabą $v$;
odwrotnie pochodna słaba musi reprezentować $D$.
Dla $\eta$ stosujemy Banacha--Steinhausa do ilorazów
różnicowych w $Y'$ jako funkcjonałów na $Y$.
Ich granica punktowa jest ograniczonym funkcjonałem
$w$, co dokładnie oznacza pochodną słabą-*.
\end{proof}

\begin{stwierdzenie}[Ciągła w normie pochodna słaba]
\label{prop:weak-to-strong-derivative14}
Niech $X,Y$ będą przestrzeniami Banacha nad $\K\in\{\R,\C\}$,
$I$ otwartym przedziałem rzeczywistym, a $\gamma:I\to X$ krzywą
mającą w każdym punkcie pochodną słabą $v(t)$.
Jeżeli $v:I\to X$ jest ciągła w normie, to $\gamma$
jest klasy $C^1$ w normie i $\gamma'=v$.
Jeżeli $\eta:I\to Y'$ ma w każdym punkcie $t$ pochodną
słabą-* $w(t)$ i $w:I\to Y'$ jest ciągła w normie,
to $\eta$ jest klasy $C^1$ w normie i $\eta'=w$.
\end{stwierdzenie}
\begin{proof}
Dla $a,b\in I$, $a<b$, ciągłość $v$ zapewnia istnienie
wektorowej całki Riemanna. Skalarne twierdzenie
Newtona--Leibniza daje dla każdego $x'\in X'$
$x'(\gamma(b)-\gamma(a))=\int_a^b x'(v(t))\,\mathrm dt$.
Ciągły funkcjonał przechodzi przez granicę sum Riemanna,
więc z rozdzielania punktów wynika
$\gamma(b)-\gamma(a)=\int_a^b v(t)\,\mathrm dt$.
Dla $h\ne0$ z $t+h\in I$, przyjmując całkę z odwróconymi
granicami ze znakiem minus, mamy
$\gamma(t+h)-\gamma(t)=\int_t^{t+h}v(s)\,\mathrm ds$.
Po odjęciu $hv(t)$, podzieleniu przez $h$ i użyciu
oszacowania normy całki otrzymujemy, że norma różnicy
ilorazu różnicowego i $v(t)$ jest co najwyżej
$\sup_{s\text{ między }t\text{ i }t+h}\|v(s)-v(t)\|$,
która dąży do zera. Ciągłość $v$ zapewnia ciągłość otrzymanej pochodnej.
W przypadku $Y'$ dla każdego $y\in Y$ stosujemy
skalarne twierdzenie Newtona--Leibniza do $\eta(\cdot)(y)$,
otrzymując $(\eta(b)-\eta(a))(y)=(\int_a^b w(s)\,\mathrm ds)(y)$.
Ewaluacje na $Y$ rozdzielają punkty $Y'$, więc
$\eta(b)-\eta(a)=\int_a^b w(s)\,\mathrm ds$.
To samo oszacowanie normy ilorazu różnicowego z $w$
daje pochodną normową i jej ciągłość.
\end{proof}

\begin{przyklad}[Pochodna słaba bez pochodnej normowej]
\label{ex:weak-not-norm-derivative14}
Niech $X=\ell^2(\N;\K)$, gdzie $\K\in\{\R,\C\}$,
$(e_n)$ jego standardowa baza ortonormalna,
oraz $\gamma:(-1,1)\to X$ krzywa określona przez $\gamma(0)=0$ i dla $0<|t|<1$ przez
$\gamma(t)=t e_{\lceil1/|t|\rceil}$. Iloraz różnicowy
w zerze jest wektorem jednostkowym, którego indeks
dąży do nieskończoności. Z twierdzenia Riesza dla każdego $x'\in X'$ istnieje
$a\in\ell^2(\N;\K)$ takie, że $x'(v)=\langle v,a\rangle$,
przy konwencji liniowości iloczynu skalarnego względem
pierwszego argumentu. Ponieważ $\sum_n|a_n|^2<\infty$,
mamy $|a_n|\to0$, a więc $x'(e_{\lceil1/|t|\rceil})\to0$.
Pochodna słaba wynosi zatem zero. Normy ilorazów są równe $1$, zatem pochodna
normowa nie istnieje. Krzywa jest ciągła w zerze,
ponieważ $\|\gamma(t)\|=|t|$. Nie twierdzimy, że jest
ciągła w pozostałych punktach.
\end{przyklad}

\begin{przyklad}[Pochodne wszystkich skalaryzacji bez pochodnej słabej]
\label{ex:scalar-not-weak-derivative14}
Niech $\K\in\{\R,\C\}$ i $c_0=c_0(\N;\K)$.
Oznaczmy przez $v_n$ ciąg mający jedynki na pierwszych
$n$ pozycjach, a zera dalej. Niech $\gamma:(-1,1)\to c_0$ będzie określona przez
$\gamma(0)=0$ i $\gamma(t)=t v_{\lceil1/|t|\rceil}$ dla $0<|t|<1$.
Identyfikujemy $(c_0)'$ z $\ell^1(\N;\K)$ przez
$a\mapsto(v\mapsto\sum_{k\ge1}a_kv_k)$. Dla każdego
$a\in\ell^1$ iloraz skalarny wynosi
$\sum_{k=1}^{\lceil1/|t|\rceil} a_k$ i dąży do
$\sum_{k\ge1}a_k$, gdyż szereg jest bezwzględnie zbieżny. Reprezentujący element $(c_0)''=\ell^\infty$
jest ciągiem samych jedynek, który nie należy do $c_0$.
Pochodna słaba w $c_0$ nie istnieje, mimo istnienia
wszystkich pochodnych skalarnych. W $\ell^\infty=(\ell^1)'$
ta sama krzywa ma natomiast pochodną słabą-* równą
$(1,1,\ldots)$ i nie ma pochodnej normowej, ponieważ
$\|v_n-(1,1,\ldots)\|_\infty=1$.
\end{przyklad}

\begin{definicja}[Różniczkowalność w sensie Bastianiego]
\label{def:bastiani14}
Niech $E,F$ będą rzeczywistymi przestrzeniami lokalnie
wypukłymi Hausdorffa i $U\subset E$ otwartym zbiorem.
Ciągłe odwzorowanie $f:U\to F$ jest klasy $C^1$ w sensie
Bastianiego, jeżeli dla wszystkich $(x,h)\in U\times E$
istnieje granica $df(x;h)=\lim_{t\to0}(f(x+th)-f(x))/t$
w topologii $F$, a $df:U\times E\to F$ jest ciągłe.
\end{definicja}
\begin{stwierdzenie}[Liniowość pochodnej Bastianiego]
\label{prop:bastiani-linear14}
Niech $E,F$ będą rzeczywistymi przestrzeniami lokalnie
wypukłymi Hausdorffa, $U\subset E$ zbiorem otwartym
i $f:U\to F$ odwzorowaniem klasy $C^1$ w sensie Bastianiego.
Wówczas dla każdego $x\in U$
odwzorowanie $h\mapsto df(x;h)$ jest liniowe i ciągłe.
\end{stwierdzenie}
\begin{proof}
Dla $\alpha\ne0$ podstawienie $u=\alpha t$ w ilorazie
różnicowym daje $df(x;\alpha h)=\alpha df(x;h)$,
a dla $\alpha=0$ iloraz jest zerowy. Dla kierunków
$h,k\in E$ i dostatecznie małego $t\ne0$ piszemy
\[
 \frac{f(x+t(h+k))-f(x)}{t}
 =\frac{f(x+th)-f(x)}t
  +\frac{f(x+th+tk)-f(x+th)}t .
\]
Dla każdego $\ell\in F'$ i ustalonego $t\ne0$ rozważmy
funkcję skalarną $g_t(u)=\ell(f(x+th+utk))$.
Z definicji pochodnej kierunkowej, ciągłości $\ell$
i wykazanej już jednorodności otrzymujemy
$g_t'(u)=\ell(df(x+th+utk;tk))=t\ell(df(x+th+utk;k))$.
Pochodna jest ciągła dzięki wspólnej ciągłości $df$.
Skalarne twierdzenie Newtona--Leibniza daje zatem
$(g_t(1)-g_t(0))/t=\int_0^1\ell(df(x+th+utk;k))\,\mathrm du$. Odcinek leży
w $U$ dla dostatecznie małego $t$, dzięki istnieniu
wypukłego otoczenia zera zawartego w $U-x$.
Dla dowolnego otoczenia zera $W$ w $E$ istnieje
wypukłe otoczenie zera $W_0\subseteq W$.
Dla małych $|t|$ oba wektory $th$ i $t(h+k)$ należą
do $W_0$, więc także ich kombinacja wypukła
$th+utk$ dla każdego $u\in[0,1]$.
Stąd $x+th+utk$ dąży do $x$ jednostajnie względem $u$
w topologii $E$. Ciągłość $df$ w $(x,k)$ oraz $\ell$
daje jednostajną zbieżność całkowanej funkcji do
$\ell(df(x;k))$. Po przejściu do granicy otrzymujemy
$\ell(df(x;h+k))=\ell(df(x;h))+\ell(df(x;k))$ dla
każdego $\ell\in F'$; twierdzenie Hahna--Banacha zapewnia
rozdzielanie punktów. Ciągłość względem $h$ jest częścią
założenia wspólnej ciągłości.
\end{proof}

\begin{uwaga}[Poziomy ogólności]
Związki między kierunkową a pełną różniczkowalnością, normą, topologią operatorową i różniczkowalnością zespoloną wykorzystują normowalność dziedziny. Wprowadzona różniczkowalność w sensie Bastianiego jest jednym z pojęć rachunku w przestrzeniach lokalnie wypukłych. Nie utożsamiamy jej z rachunkiem Frécheta bez dodatkowych założeń. Pochodne krzywych bezwzględnie ciągłych i całek słabych są rozwijane w \hyperref[sec:integral-vector15b]{module 15B}.
\end{uwaga}

% Wersja źródłowa: audyt II, jednostka 15

\section{Pochodne wyższych rzędów i lokalne twierdzenia analizy nieliniowej}
\label{sec:wyzsze-pochodne15}

\providecommand{\atlasref}[2]{#2}


\subsection{Odwzorowania wieloliniowe i definicja wyższych pochodnych}

\noindent\emph{Powiązanie tematyczne w Bibliotece matematycznej:} \atlasref{LA-08::def:lin8-wieloliniowe}{LA-08: odwzorowania wieloliniowe}.


W podstawowym rachunku różniczkowym różniczkę $Df(x)$ traktowaliśmy jako ograniczony operator liniowy. Przy dalszym różniczkowaniu otrzymujemy odwzorowania wieloliniowe. Ta perspektywa pozwala sformułować twierdzenie Taylora i twierdzenia o funkcji odwrotnej oraz uwikłanej bez redukowania przestrzeni Banacha do współrzędnych.

\begin{definicja}[Ciągłe odwzorowania wieloliniowe]
\label{def:multilinear15}
Niech $m\ge1$, $\K\in\{\R,\C\}$, a $E_1,\ldots,E_m,F$ będą przestrzeniami unormowanymi nad $\K$. Przez $\mathcal L^m(E_1,\ldots,E_m;F)$ oznaczamy przestrzeń odwzorowań $A:E_1\times\cdots\times E_m\to F$, liniowych w każdym argumencie i ciągłych, wyposażoną w normę
$\|A\|=\sup\{\|A(h_1,\ldots,h_m)\|:\|h_j\|\le1,\ 1\le j\le m\}$. Jeśli wszystkie dziedziny są równe $E$, piszemy $\mathcal L^m(E;F)$, a dla $m=1$ używamy dotychczasowego oznaczenia $\mathcal L(E,F)$. Przez $\mathcal L^m_s(E;F)$ oznaczamy podprzestrzeń odwzorowań symetrycznych, tj. niezmienniczych przy permutacjach argumentów.
\end{definicja}

\begin{stwierdzenie}[Kryterium ograniczoności i zupełność]
\label{prop:multi-banach15}
Niech $m\ge1$, $\K\in\{\R,\C\}$, $E_1,\ldots,E_m$ będą przestrzeniami unormowanymi nad $\K$, a $F$ przestrzenią Banacha nad tym samym ciałem. Odwzorowanie $m$-liniowe $A:E_1\times\cdots\times E_m\to F$ jest ciągłe wtedy i tylko wtedy, gdy istnieje $C\ge0$ takie, że
$\|A(h_1,\ldots,h_m)\|\le C\prod_{j=1}^m\|h_j\|$ dla wszystkich $h_j\in E_j$. Przestrzeń $\mathcal L^m(E_1,\ldots,E_m;F)$ jest przestrzenią Banacha.
\end{stwierdzenie}
\begin{proof}
Jeśli $A$ jest ciągłe w zerze, istnieje $\delta>0$ takie, że $\|A(v_1,\ldots,v_m)\|\le1$ dla $\|v_j\|<\delta$. Dla niezerowych $h_j$ podstawiamy $v_j=\delta h_j/(2\|h_j\|)$, co po wykorzystaniu wieloliniowości daje oszacowanie z $C=(2/\delta)^m$. Gdy któryś argument jest zerowy, nierówność jest oczywista. Odwrotnie, ustalmy $(h_1,\ldots,h_m)$. Po rozwinięciu wieloliniowym różnica $A(h_1+k_1,\ldots,h_m+k_m)-A(h_1,\ldots,h_m)$ jest sumą wyrazów indeksowanych przez niepuste podzbiory $J\subseteq\{1,\ldots,m\}$. Norma wyrazu o indeksie $J$ nie przekracza $C\prod_{j\in J}\|k_j\|\prod_{j\notin J}\|h_j\|$. Każdy z tych skończenie wielu wyrazów dąży do zera, gdy $\max_j\|k_j\|\to0$. Otrzymujemy ciągłość w każdym punkcie.

Niech $(A_n)$ będzie ciągiem Cauchy'ego w normie operatorowej. Dla ustalonych $h_1,\ldots,h_m$ ciąg $A_n(h_1,\ldots,h_m)$ jest Cauchy'ego w $F$, więc definiuje $A(h_1,\ldots,h_m)$. Granica zachowuje wieloliniowość. Ponieważ normy $\|A_n\|$ są ograniczone przez wspólną stałą, $A$ spełnia oszacowanie wieloliniowe. Przechodząc do granicy w nierówności $\|(A_n-A_p)(h_1,\ldots,h_m)\|\le\varepsilon\prod\|h_j\|$ dla $n,p\ge N$, otrzymujemy $\|A_n-A\|\le\varepsilon$. To dowodzi zupełności.
\end{proof}

\begin{lemat}[Utożsamienie operatorów wieloliniowych]
\label{lem:currying15}
Niech $m\ge1$, $\K\in\{\R,\C\}$, a $E,F$ będą przestrzeniami unormowanymi nad $\K$. Naturalne odwzorowanie
$\mathcal L(E,\mathcal L^m(E;F))\to\mathcal L^{m+1}(E;F)$, zadane przez $T\mapsto\widetilde T$ oraz $\widetilde T(h_0,h_1,\ldots,h_m)=(Th_0)(h_1,\ldots,h_m)$, jest izometrycznym izomorfizmem liniowym.
\end{lemat}
\begin{proof}
Liniowość, wieloliniowość i ciągłość wynikają z oszacowania $\|\widetilde T(h_0,\ldots,h_m)\|\le\|T\|\prod_{j=0}^m\|h_j\|$. Ponadto, z definicji norm operatorowych,
\[
\begin{aligned}
\|\widetilde T\|&=\sup_{\|h_0\|\le1}\sup_{\substack{\|h_j\|\le1\\1\le j\le m}}\|(Th_0)(h_1,\ldots,h_m)\|\\
&=\|T\|.
\end{aligned}
\]
Dla $A\in\mathcal L^{m+1}(E;F)$ definiujemy $(T_Ah_0)(h_1,\ldots,h_m)=A(h_0,h_1,\ldots,h_m)$. Wieloliniowość $A$ daje liniowość $T_A$, zaś $\|T_Ah_0\|\le\|A\|\|h_0\|$ zapewnia jego ograniczoność. Ponieważ $\widetilde T_A=A$, otrzymujemy surjektywność, a izometryczność daje także iniektywność.
\end{proof}

\begin{definicja}[Klasa $C^m$ i pochodne wyższych rzędów]
\label{def:ck15}
Niech $m\ge1$, $\K\in\{\R,\C\}$, $E,F$ będą przestrzeniami Banacha nad $\K$, a $U\subseteq E$ zbiorem otwartym. Wszystkie pochodne Frécheta rozpatrujemy nad $\K$ (dla $\K=\C$ są zespolenie liniowe). Mówimy, że $f:U\to F$ jest klasy $C^m$, gdy $f\in C^1$ oraz rekurencyjnie $D^{j-1}f:U\to\mathcal L^{j-1}(E;F)$ jest klasy $C^1$ dla $2\le j\le m$. Pochodną $D^jf(x)$ utożsamiamy przez Lemat~\ref{lem:currying15} z elementem $\mathcal L^j(E;F)$. Przyjmujemy $D^0f=f$ oraz $C^\infty(U,F)=\bigcap_{m\ge1}C^m(U,F)$. Dla $1\le j\le m$ ciągłość $D^jf$ rozumiemy w normie $\mathcal L^j(E;F)$.
\end{definicja}

\begin{przyklad}[Odwzorowanie wieloliniowe i wielomiany]
\label{ex:multi-poly15}
Niech $E,F$ będą rzeczywistymi przestrzeniami Banacha, $m\ge1$, $B\in\mathcal L^m(E;F)$ i $P:E\to F$ będzie dane wzorem $P(x)=B(x,\ldots,x)$. Wówczas $P\in C^\infty(E,F)$ oraz
$DP(x)h=\sum_{j=1}^m B(x,\ldots,x,\underset{j}{h},x,\ldots,x)$.
Istotnie, różnica $P(x+h)-P(x)$ składa się z wyrazów, w których $h$ występuje co najmniej raz; wyrazy z dokładnie jednym $h$ tworzą wskazany operator, a suma pozostałych jest $O(\|h\|^2)$ przy ustalonym $x$. Indukcyjnie pochodna rzędu $j\le m$ jest sumą po wszystkich iniekcjach $\{1,\ldots,j\}\to\{1,\ldots,m\}$: w pozycjach wskazanych przez iniekcję umieszczamy kolejno $h_1,\ldots,h_j$, a w pozostałych $x$. Różniczkując w następnym kierunku, zastępujemy jeden z pozostałych czynników $x$ przez $h_{j+1}$. Wyrazy z co najmniej dwoma nowymi przyrostami dają resztę $O(\|h_{j+1}\|^2)$ jednostajnie dla $\|h_i\|\le1$, z oszacowania wieloliniowego. To uzasadnia indukcję w normie operatorowej. Ciągłość każdej pochodnej wynika z tego samego oszacowania. Pochodna rzędu $m$ jest stała, a pochodne od rzędu $m+1$ są zerowe. Dla symetrycznego $B$ zachodzi $DP(x)h=mB(x,\ldots,x,h)$, a $D^mP(x)=m!B$.
\end{przyklad}

\subsection{Symetria drugiej pochodnej i reguły rachunku}

\begin{twierdzenie}[Symetria drugiej pochodnej]
\label{thm:schwarz15}
Niech $E,F$ będą rzeczywistymi przestrzeniami Banacha, $U\subseteq E$ otwarty i $f\in C^2(U,F)$. Dla każdego $x\in U$ i $h,k\in E$ zachodzi
$D^2f(x)[h,k]=D^2f(x)[k,h]$. W szczególności $D^2f(x)\in\mathcal L_s^2(E;F)$.
\end{twierdzenie}
\begin{proof}
Dla małych rzeczywistych $s,t$ połóżmy
$\Delta(s,t)=f(x+sh+tk)-f(x+sh)-f(x+tk)+f(x)$. Z dwukrotnego zastosowania wzoru całkowego na odcinku dla funkcji klasy $C^1$ otrzymujemy
\[
\Delta(s,t)=st\int_0^1\int_0^1 D^2f(x+\alpha sh+\beta tk)[k,h] \,\mathrm d\beta\,\mathrm d\alpha.
\]
Istotnie, najpierw piszemy $\Delta(s,t)=s\int_0^1(Df(x+\alpha sh+tk)-Df(x+\alpha sh))h\,\mathrm d\alpha$, a następnie stosujemy wzór na przyrost do $Df$ w kierunku $tk$ i ewaluujemy na $h$. Zgodnie z utożsamieniem z Lematu~\ref{lem:currying15} pierwszy argument $D^2f$ jest kierunkiem różniczkowania $Df$. Jeśli zaczniemy od kierunku $k$, otrzymamy natomiast $st\int_0^1\int_0^1D^2f(x+\alpha sh+\beta tk)[h,k]\,\mathrm d\alpha\,\mathrm d\beta$. Nie trzeba zakładać symetrii ani zamieniać kolejności całkowania, aby porównać granice obu wyrażeń. Dla $s,t\ne0$ dzielimy przez $st$ i przechodzimy z $(s,t)$ do $(0,0)$. Ciągłość $D^2f$ w normie operatorowej zapewnia jednostajną zbieżność funkcji podcałkowych na kwadracie $[0,1]^2$ odpowiednio do $D^2f(x)[k,h]$ i $D^2f(x)[h,k]$. Obie granice są równe, co kończy dowód.
\end{proof}


\begin{wniosek}[Symetria wszystkich pochodnych wyższych rzędów]
\label{cor:symmetric-higher15}
Niech $E,F$ będą rzeczywistymi przestrzeniami Banacha, $U\subseteq E$ otwarty, $m\ge2$ i $f\in C^m(U,F)$. Wówczas $D^mf(x)\in\mathcal L_s^m(E;F)$ dla każdego $x\in U$.
\end{wniosek}
\begin{proof}
Stosujemy indukcję po $m$. Dla $m=2$ jest to Twierdzenie~\ref{thm:schwarz15}. Jeśli $m\ge3$, to odwzorowanie $D^{m-2}f:U\to\mathcal L^{m-2}(E;F)$ jest klasy $C^2$, a więc zastosowanie tego twierdzenia do $D^{m-2}f$ daje symetrię $D^mf(x)$ względem pierwszych dwóch argumentów. Z założenia indukcyjnego $D^{m-1}f(z)$ jest symetryczne dla każdego $z\in U$. Różniczkując tożsamości wyrażające tę symetrię w kierunku pierwszego argumentu, stwierdzamy, że $D^mf(x)$ jest symetryczne także względem ostatnich $m-1$ argumentów. Łącznie otrzymujemy niezmienniczość przy wszystkich zamianach sąsiednich argumentów, a te generują wszystkie permutacje.
\end{proof}

\begin{stwierdzenie}[Druga pochodna złożenia]
\label{prop:second-chain15}
Niech $E,F,G$ będą rzeczywistymi przestrzeniami Banacha, $U\subseteq E$ i $V\subseteq F$ otwarte, $f\in C^2(U,F)$, $g\in C^2(V,G)$ oraz $f(U)\subseteq V$. Wówczas $g\circ f\in C^2(U,G)$ i dla $x\in U$ oraz $h,k\in E$ zachodzi
\[
\begin{aligned}
D^2(g\circ f)(x)[h,k]
={}&D^2g(f(x))[Df(x)h,Df(x)k]\\
 &+Dg(f(x))\bigl(D^2f(x)[h,k]\bigr).
\end{aligned}
\]
\end{stwierdzenie}
\begin{proof}
Z reguły łańcuchowej $D(g\circ f)(x)=Dg(f(x))\circ Df(x)$. Kompozycja operatorów jest ciągłym odwzorowaniem dwuliniowym $\mathcal L(F,G)\times\mathcal L(E,F)\to\mathcal L(E,G)$. Stosując regułę różniczkowania ciągłego odwzorowania dwuliniowego do odwzorowań $x\mapsto Dg(f(x))$ i $x\mapsto Df(x)$, otrzymujemy po zadziałaniu na $k$ dwa składniki. Pierwszy, dzięki regule łańcuchowej zastosowanej do $Dg$, ma postać $D^2g(f(x))[Df(x)h,Df(x)k]$, a drugi jest równy $Dg(f(x))(D^2f(x)[h,k])$. Ciągłość wszystkich występujących pochodnych dowodzi przynależności złożenia do $C^2$.
\end{proof}

\begin{przyklad}[Hesjan funkcjonału na przestrzeni Hilberta]
\label{ex:hessian15}
Niech $H$ będzie rzeczywistą przestrzenią Hilberta, $A\in\mathcal L(H)$ i $b\in H$. Dla $J(x)=\frac12\langle Ax,x\rangle-\langle x,b\rangle$ mamy
$DJ(x)h=\langle h,(A+A^*)x/2-b\rangle$ oraz
$D^2J(x)[h,k]=\langle h,(A+A^*)k/2\rangle$.
Przy konwencji liniowości iloczynu skalarnego w pierwszym argumencie wzór jest symetryczny dla rzeczywistego $H$, ponieważ $(A+A^*)/2$ jest samosprzężony. Jeśli $A=A^*$, to $D^2J(x)[h,k]=\langle h,Ak\rangle$.
\end{przyklad}

\subsection{Twierdzenie Taylora w przestrzeniach Banacha}

\noindent\emph{Powiązanie tematyczne w Bibliotece matematycznej:} \atlasref{AM1-23::ma1:thm:taylor-peano}{AM1-23: twierdzenie Taylora z resztą Peana}.


\begin{twierdzenie}[Wzór Taylora z resztą całkową]
\label{thm:taylor15}
Niech $E,F$ będą rzeczywistymi przestrzeniami Banacha, $U\subseteq E$ otwarty, $m\ge1$, $f\in C^m(U,F)$, a $x\in U$ i $h\in E$ takie, że odcinek $\{x+th:0\le t\le1\}$ zawiera się w $U$. Wtedy
\[
 f(x+h)=\sum_{j=0}^{m-1}\frac1{j!}D^jf(x)[h^j]
 +\frac1{(m-1)!}\int_0^1(1-t)^{m-1}D^mf(x+th)[h^m] \,\mathrm dt,
\]
gdzie $D^0f(x)[h^0]=f(x)$, a $D^jf(x)[h^j]$ oznacza $D^jf(x)[h,\ldots,h]$. Ponadto
$f(x+h)=\sum_{j=0}^{m}\frac1{j!}D^jf(x)[h^j]+r_m(h)$, gdzie
\[
\|r_m(h)\|\le \frac{\|h\|^m}{m!}\sup_{0\le t\le1}\|D^mf(x+th)-D^mf(x)\|.
\]
W szczególności $\|r_m(h)\|=o(\|h\|^m)$ dla $h\to0$.
\end{twierdzenie}
\begin{proof}
Ustalmy $x,h$ i zdefiniujmy $\gamma(t)=f(x+th)$. Z reguły łańcuchowej i indukcji mamy $\gamma^{(j)}(t)=D^jf(x+th)[h^j]$ dla $1\le j\le m$. Pierwsza pochodna jest ciągła, zatem stosujemy podstawowe twierdzenie rachunku całkowego dla funkcji o wartościach w $F$. Udowodnimy wzór dla krzywej $\gamma$ przez indukcję po $m$. Dla $m=1$ jest to twierdzenie Newtona--Leibniza. W kroku od $m$ do $m+1$ zakładamy $\gamma\in C^{m+1}$ i w reszcie rzędu $m$ całkujemy przez części: dla każdego $\ell\in F'$ stosujemy skalarny wzór do $\ell(\gamma^{(m)}(t))$ i funkcji $-(1-t)^m/m!$, której pochodna wynosi $(1-t)^{m-1}/(m-1)!$. Otrzymujemy
\[
 \frac1{(m-1)!}\int_0^1(1-t)^{m-1}\gamma^{(m)}(t)\,\mathrm dt
 =\frac{\gamma^{(m)}(0)}{m!}
 +\frac1{m!}\int_0^1(1-t)^m\gamma^{(m+1)}(t)\,\mathrm dt .
\]
Równość zachodzi w $F$, ponieważ ciągłe funkcjonały przechodzą przez całkę Riemanna i rozdzielają punkty. Podstawienie do wzoru indukcyjnego daje przypadek $m+1$. Wreszcie wstawiamy wyrażenia na $\gamma^{(j)}$.

Odejmujemy od wyrażenia pod całką stałą $D^mf(x)[h^m]$, której całka z wagą $(1-t)^{m-1}/(m-1)!$ jest równa $D^mf(x)[h^m]/m!$. Następnie stosujemy normę operatorową, oszacowanie normy całki przez całkę normy oraz $\int_0^1(1-t)^{m-1}\,\mathrm dt=1/m$. Ciągłość $D^mf$ w $x$ zapewnia, że supremum w oszacowaniu reszty dąży do zera przy $h\to0$.
\end{proof}

\begin{wniosek}[Wzór drugiego rzędu]
\label{cor:taylor2-15}
Niech $E,F$ będą rzeczywistymi przestrzeniami Banacha, $U\subseteq E$ zbiorem otwartym, $f\in C^2(U,F)$ oraz $x\in U$, $h\in E$ takie, że $[x,x+h]\subseteq U$. Wówczas zachodzi
\[
 f(x+h)=f(x)+Df(x)h+\tfrac12D^2f(x)[h,h]
 +\int_0^1(1-t)\bigl(D^2f(x+th)-D^2f(x)\bigr)[h,h] \,\mathrm dt.
\]
Jeżeli $M\ge0$ i dla każdego $z\in[x,x+h]$ zachodzi $\|D^2f(z)\|\le M$, to reszta po wyrazach rzędu pierwszego ma normę co najwyżej $M\|h\|^2/2$.
\end{wniosek}

\begin{stwierdzenie}[Warunki konieczne i wystarczające minimum]
\label{prop:extremum15}
Niech $E$ będzie rzeczywistą przestrzenią Banacha, $U\subseteq E$ otwarty, $x_0\in U$, a $f\in C^2(U,\R)$. Jeżeli $x_0$ jest lokalnym minimum $f$, to $Df(x_0)=0$ i $D^2f(x_0)[h,h]\ge0$ dla każdego $h\in E$. Odwrotnie, jeżeli $Df(x_0)=0$ oraz istnieje $c>0$ takie, że $D^2f(x_0)[h,h]\ge c\|h\|^2$ dla wszystkich $h\in E$, to $x_0$ jest ścisłym lokalnym minimum $f$; dokładniej, dla dostatecznie małych $h$ zachodzi $f(x_0+h)\ge f(x_0)+(c/4)\|h\|^2$.
\end{stwierdzenie}
\begin{proof}
Dla dowolnego $h$ rozważamy rzeczywistą funkcję $\varphi(t)=f(x_0+th)$. Ma ona minimum w $0$, więc $\varphi'(0)=Df(x_0)h=0$ i $\varphi''(0)=D^2f(x_0)[h,h]\ge0$. W drugim kierunku ciągłość $D^2f$ daje promień $r>0$ taki, że $\|D^2f(x_0+u)-D^2f(x_0)\|\le c/2$ dla $\|u\|<r$. Stąd $D^2f(x_0+u)[h,h]\ge(c/2)\|h\|^2$. Z całkowego wzoru Taylora drugiego rzędu, przy $Df(x_0)=0$, otrzymujemy $f(x_0+h)-f(x_0)\ge(c/2)\|h\|^2\int_0^1(1-t)\,\mathrm dt=c\|h\|^2/4$ dla $\|h\|<r$.
\end{proof}

\begin{przyklad}[Dodatniość bez jednostajności nie wystarcza]
\label{ex:positive-hessian15}
W rzeczywistej przestrzeni Hilberta $H=\ell^2(\N;\R)$ niech
$Q(x)=\frac12\sum_{n=1}^\infty x_n^2/n$, zaś $f(x)=Q(x)-\|x\|_2^3$. Operator diagonalny $A:H\to H$, $(Ax)_n=x_n/n$, jest samosprzężony i ograniczony, ponieważ $\|Ax\|_2^2\le\|x\|_2^2$. Zatem $Q(x)=\langle Ax,x\rangle/2$ jest gładką funkcją kwadratową. Poniżej uzasadniamy, że $G(x)=\|x\|_2^3$ jest klasy $C^2$, a więc również $f$ ma tę regularność. Mamy $Df(0)=0$ i $D^2f(0)[h,h]=\sum_{n\ge1}h_n^2/n>0$ dla $h\ne0$, ale $0$ nie jest minimum lokalnym. Istotnie, dla $t_n=1/n$ mamy $t_ne_n\to0$ w normie i $f(t_ne_n)=1/(2n^3)-1/n^3=-1/(2n^3)<0=f(0)$. Funkcja $G(x)=\|x\|_2^3$ jest klasy $C^2$ również w zerze. Dla $x\ne0$ mamy
$DG(x)h=3\|x\|\langle x,h\rangle$ oraz
$D^2G(x)[h,k]=3\|x\|\langle h,k\rangle+3\langle x,h\rangle\langle x,k\rangle/\|x\|$.
Stąd $\|D^2G(x)\|\le6\|x\|$. Ponieważ $|G(h)-G(0)|/\|h\|=\|h\|^2\to0$, mamy $DG(0)=0$. Ponadto $\|DG(x)\|\le3\|x\|^2$, co dowodzi $D^2G(0)=0$ i ciągłości drugiej pochodnej w zerze.
\end{przyklad}

\subsection{Szereg Neumanna i odwracanie operatorów}

\begin{lemat}[Szereg Neumanna]
\label{lem:neumann15}
Niech $E$ będzie przestrzenią Banacha nad $\K\in\{\R,\C\}$, $I$ operatorem identycznościowym i $T\in\mathcal L(E)$ operatorem liniowym ograniczonym spełniającym $\|T\|<1$. Wtedy $I-T$ jest odwracalny oraz
$(I-T)^{-1}=\sum_{n=0}^\infty T^n$ w normie operatorowej, przy czym $\|(I-T)^{-1}\|\le(1-\|T\|)^{-1}$.
\end{lemat}
\begin{proof}
Szereg jest absolutnie zbieżny w przestrzeni Banacha $\mathcal L(E)$. Dla sum częściowych $S_N=\sum_{n=0}^N T^n$ zachodzi $(I-T)S_N=S_N(I-T)=I-T^{N+1}$; przechodząc do granicy, otrzymujemy obustronną odwrotność. Oszacowanie normy wynika z szeregu geometrycznego.
\end{proof}

\begin{stwierdzenie}[Otwartość grupy operatorów odwracalnych]
\label{prop:inverse-operator15}
Niech $E$ będzie rzeczywistą przestrzenią Banacha i $\mathrm{GL}(E)$ zbiorem odwracalnych operatorów w $\mathcal L(E)$. Wówczas $\mathrm{GL}(E)$ jest otwarty, a $\iota:\mathrm{GL}(E)\to\mathcal L(E)$, $\iota(A)=A^{-1}$, jest klasy $C^\infty$. Jej pochodna ma postać
$D\iota(A)H=-A^{-1}HA^{-1}$.
\end{stwierdzenie}
\begin{proof}
Dla $A\in\mathrm{GL}(E)$ i $H\in\mathcal L(E)$ spełniającego $\|A^{-1}H\|<1$ mamy $A+H=A(I+A^{-1}H)$, więc Lemat~\ref{lem:neumann15} daje odwracalność oraz lokalną jednostajną ograniczoność odwrotności. Z tożsamości
$(A+H)^{-1}-A^{-1}=-(A+H)^{-1}HA^{-1}$ wynika
\[
(A+H)^{-1}-A^{-1}+A^{-1}HA^{-1}
=A^{-1}H(A+H)^{-1}HA^{-1}.
\]
Dla $\|A^{-1}H\|\le1/2$ szereg Neumanna daje $\|(A+H)^{-1}\|\le2\|A^{-1}\|$. Norma prawej strony nie przekracza zatem $2\|A^{-1}\|^3\|H\|^2$, co dowodzi istnienia wskazanej pochodnej Frécheta. Jej ciągłość otrzymujemy z ciągłości odwracania i ciągłości kompozycji operatorów. Powtarzając różniczkowanie wzoru $D\iota(A)H=-A^{-1}HA^{-1}$ z regułą wieloliniową, otrzymujemy rekurencyjnie pochodne wszystkich rzędów i ich ciągłość.
\end{proof}

\subsection{Twierdzenie o funkcji odwrotnej}

\noindent\emph{Powiązanie tematyczne w Bibliotece matematycznej:} \atlasref{DG-13::thm:wspolrzedne-normalne}{DG-13: współrzędne normalne jako zastosowanie w skończonym wymiarze}.

\begin{twierdzenie}[Twierdzenie o funkcji odwrotnej dla przestrzeni Banacha]
\label{thm:inverse15}
Niech $E,F$ będą rzeczywistymi przestrzeniami Banacha, $U\subseteq E$ otwarty, $f\in C^1(U,F)$ i $x_0\in U$. Załóżmy, że $A=Df(x_0):E\to F$ jest izomorfizmem liniowym i topologicznym, tj. ma ograniczoną odwrotność $A^{-1}:F\to E$. Wówczas istnieją otoczenia otwarte $U_0\subseteq U$ punktu $x_0$ i $V_0\subseteq F$ punktu $y_0=f(x_0)$ takie, że $f|_{U_0}:U_0\to V_0$ jest bijekcją, a odwrotność $g:V_0\to U_0$ należy do $C^1(V_0,E)$ i spełnia
$Dg(y)=(Df(g(y)))^{-1}$ dla każdego $y\in V_0$.
\end{twierdzenie}
\begin{proof}
\emph{Krok 1: jednostajne przybliżenie liniowe.} Z ciągłości $Df$ w $x_0$ wybieramy $r>0$ tak małe, aby domknięta kula $\overline B(x_0,r)$ zawierała się w $U$ i
$\|I-A^{-1}Df(x)\|\le q<1$ dla $x\in\overline B(x_0,r)$; można przyjąć $q=1/2$. Dla $y\in F$ określamy $T_y(x)=x-A^{-1}(f(x)-y)$. Dla $x,z$ w tej kuli wzór całkowy na odcinku daje
\[
\|T_y(x)-T_y(z)\|\le q\|x-z\|.
\]
Ponadto $T_y(x_0)-x_0=A^{-1}(y-y_0)$. Dla $\|A^{-1}(y-y_0)\|\le(1-q)r$ operator $T_y$ przeprowadza $\overline B(x_0,r)$ w siebie, bo
$\|T_y(x)-x_0\|\le q\|x-x_0\|+\|A^{-1}(y-y_0)\|\le r$.

\emph{Krok 2: punkt stały.} Przypomnijmy zasadę kontrakcji: jeżeli $T$ jest odwzorowaniem $q$-lipschitzowskim pełnej przestrzeni metrycznej w siebie, $q<1$, to ciąg $x_{n+1}=T(x_n)$ jest Cauchy'ego, ponieważ $d(x_{n+1},x_n)\le q^n d(x_1,x_0)$. Dla $p>n$ nierówność trójkąta daje $d(x_p,x_n)\le q^n d(x_1,x_0)/(1-q)$, więc z zupełności ciąg ma granicę $x_*$. Ciągłość $T$ daje $T(x_*)=\lim_n T(x_n)=\lim_n x_{n+1}=x_*$. Dwa punkty stałe $u,v$ spełniają $d(u,v)\le qd(u,v)$, zatem są równe. Stosujemy tę zasadę do $T_y$ na domkniętej kuli w $E$. Otrzymujemy jedyny $x$ z $T_y(x)=x$, co jest równoważne $f(x)=y$. Jeżeli $\|A^{-1}(y-y_0)\|<(1-q)r$, to dla punktu stałego $x$ zachodzi $(1-q)\|x-x_0\|\le\|A^{-1}(y-y_0)\|<(1-q)r$, więc $x$ leży we wnętrzu kuli. Dodatkowo każdy $x\in B(x_0,r)$ ma na tym obszarze co najwyżej jeden punkt o tej samej wartości $f(x)$: jeżeli $f(x)=f(z)$, to $x,z$ są punktami stałymi tego samego odwzorowania $T_{f(x)}$, dla którego oszacowanie Lipschitza daje $\|x-z\|\le q\|x-z\|$, a więc $x=z$.

\emph{Krok 3: otwarte otoczenia i ciągłość odwrotności.} Ustalmy $0<r_0<r$ i połóżmy $U_0=B(x_0,r_0)$ oraz $V_0=f(U_0)$. Z oszacowania
$\|(x-z)-A^{-1}(f(x)-f(z))\|\le q\|x-z\|$ wynika
\[
\|x-z\|\le\frac{\|A^{-1}\|}{1-q}\,\|f(x)-f(z)\| \qquad(x,z\in B(x_0,r)).
\]
Zatem $f|_{U_0}$ jest różnowartościowe, a odwrotność $g:V_0\to U_0$ jest Lipschitzowska ze stałą $\|A^{-1}\|/(1-q)$. Zbiór $V_0$ jest otwarty: dla dowolnego $x\in U_0$ wybieramy kulę $\overline B(x,\delta)\subset U_0$; dla $\|A^{-1}(y-f(x))\|<(1-q)\delta$ i $z\in\overline B(x,\delta)$ zachodzi $\|T_y(z)-x\|\le q\|z-x\|+\|A^{-1}(y-f(x))\|<\delta$. Zasada kontrakcji daje rozwiązanie $f(z)=y$ w tej kuli. Warunek na $y$ określa otwarte otoczenie $f(x)$. W konsekwencji otoczenie $f(x)$ zawiera się w $V_0$.

\emph{Krok 4: różniczkowalność odwrotności.} Z wyboru $r$ mamy $\|I-A^{-1}Df(x)\|\le q<1$ dla $x\in U_0$. Lemat~\ref{lem:neumann15} daje odwracalność $A^{-1}Df(x)$ w $\mathcal L(E)$, a z $Df(x)=A(A^{-1}Df(x))$ wynika odwracalność $Df(x):E\to F$. Ustalmy $y=f(x)\in V_0$ i rozważmy $y+k=f(x+h)$, gdzie $h=g(y+k)-g(y)\to0$. Z różniczkowalności $f$ w $x$:
$k=Df(x)h+\rho(h)$ z $\|\rho(h)\|=o(\|h\|)$. Po pomnożeniu przez $(Df(x))^{-1}$ dostajemy
$h-(Df(x))^{-1}k=-(Df(x))^{-1}\rho(h)=o(\|h\|)=o(\|k\|)$,
gdyż otrzymana w kroku 3 nierówność zapewnia $\|h\|\le C\|k\|$. Stąd $Dg(y)=(Df(x))^{-1}$. Aby uzasadnić ciągłość odwracania także dla operatorów między różnymi przestrzeniami, oznaczmy przez $C=Df(x)$ i $B=Df(z)$ dla $z$ bliskich $x$. Mamy $B^{-1}=(I+C^{-1}(B-C))^{-1}C^{-1}$, a szereg Neumanna daje $B^{-1}\to C^{-1}$ przy $B\to C$ w normie operatorowej. Ciągłość $Df$ i $g$ daje więc ciągłość $Dg$. Tak więc $g\in C^1$.
\end{proof}

\begin{wniosek}[Wersja $C^m$]
\label{cor:inverse-ck15}
Niech $E,F$ będą rzeczywistymi przestrzeniami Banacha, $U\subseteq E$ otwarty, $x_0\in U$, a $f:U\to F$ odwzorowaniem klasy $C^m$, gdzie $m\ge2$, albo klasy $C^\infty$. Jeżeli $Df(x_0):E\to F$ jest izomorfizmem liniowym i topologicznym, to lokalna odwrotność $g$ z Twierdzenia~\ref{thm:inverse15} ma tę samą regularność co $f$.
\end{wniosek}
\begin{proof}
Ustalmy $A=Df(x_0)$. Wzór $Dg(y)=(A^{-1}Df(g(y)))^{-1}A^{-1}$ sprowadza odwracanie do $\mathrm{GL}(E)$, gdzie Stwierdzenie~\ref{prop:inverse-operator15} zapewnia gładkość. Reguła łańcuchowa wyższych rzędów w potrzebnej tu postaci wynika indukcyjnie z reguły pierwszego rzędu: pochodna złożenia jest iloczynem dwóch odwzorowań operatorowych, a ciągłe mnożenie dwuliniowe zachowuje klasę $C^j$ przez wielokrotne stosowanie reguły iloczynu. Możemy zatem prowadzić indukcję: jeśli $g\in C^j$ i $f\in C^{j+1}$, to $Df\circ g\in C^j$, zatem także $Dg\in C^j$, a więc $g\in C^{j+1}$. Zaczynamy od $g\in C^1$.
\end{proof}

\begin{przyklad}[Nieliniowe odwzorowanie na przestrzeni Hilberta]
\label{ex:inverse-hilbert15}
Niech $H$ będzie rzeczywistą przestrzenią Hilberta, $v\in H$ ustalony i $f:H\to H$ dane przez $f(x)=x+\|x\|^2v$. Wtedy $Df(x)h=h+2\langle x,h\rangle v$, a $Df(0)=I$. Twierdzenie~\ref{thm:inverse15} zapewnia lokalną odwrotność w otoczeniu zera. Jeżeli $v\ne0$, odwzorowanie nie jest nawet iniektywne na całym $H$: $f(0)=0$ i $f(-v/\|v\|^2)=0$, choć $-v/\|v\|^2\ne0$.
\end{przyklad}

\subsection{Twierdzenie o funkcji uwikłanej}

\begin{twierdzenie}[Twierdzenie o funkcji uwikłanej dla przestrzeni Banacha]
\label{thm:implicit15}
Niech $X,Y,Z$ będą rzeczywistymi przestrzeniami Banacha, $\mathcal U\subseteq X\times Y$ zbiorem otwartym, $F\in C^1(\mathcal U,Z)$, $(x_0,y_0)\in\mathcal U$ i $F(x_0,y_0)=0$. Załóżmy, że pochodna częściowa $D_yF(x_0,y_0):Y\to Z$ jest izomorfizmem liniowym i topologicznym. Wtedy istnieją otwarte otoczenia $U_0\subseteq X$ punktu $x_0$, $V_0\subseteq Y$ punktu $y_0$ takie, że $U_0\times V_0\subseteq\mathcal U$, oraz jedyna funkcja $\varphi:U_0\to V_0$ klasy $C^1$ taka, że dla $(x,y)\in U_0\times V_0$ zachodzi
$F(x,y)=0\ \Longleftrightarrow\ y=\varphi(x)$. Ponadto
\[
D\varphi(x)=-\bigl(D_yF(x,\varphi(x))\bigr)^{-1}D_xF(x,\varphi(x))\qquad(x\in U_0).
\]
Jeżeli $F\in C^m$ dla $m\ge2$ lub $C^\infty$, to również $\varphi$ ma odpowiednią regularność.
\end{twierdzenie}
\begin{proof}
Definiujemy $\Psi:\mathcal U\to X\times Z$ wzorem $\Psi(x,y)=(x,F(x,y))$. Jest to odwzorowanie klasy $C^1$ i
$D\Psi(x_0,y_0)(h,k)=(h,Bh+Ak)$, gdzie $A=D_yF(x_0,y_0)$ i $B=D_xF(x_0,y_0)$. Operator ten jest odwracalny: odwrotny przeprowadza $(h,z)$ w $(h,A^{-1}(z-Bh))$ i jest ograniczony. Twierdzenie~\ref{thm:inverse15} daje lokalny dyfeomorfizm $\Psi$ między otoczeniami $(x_0,y_0)$ oraz $(x_0,0)$. Oznaczmy dziedzinę i obraz lokalnego dyfeomorfizmu przez $\mathcal O$ i $\mathcal P$. Wybieramy otwarte otoczenia $U_1,V_0,W_0$ odpowiednio punktów $x_0,y_0,0$ takie, że $U_1\times V_0\subseteq\mathcal O$ i $U_1\times W_0\subseteq\mathcal P$; istnienie wynika z topologii produktu i otwartości $\mathcal O,\mathcal P$. Zapiszmy odwrotność $\Psi^{-1}(x,z)=(x,G(x,z))$; pierwsza współrzędna ma tę postać, ponieważ $\Psi$ zachowuje $x$. Ponieważ $G(x_0,0)=y_0$ i $G$ jest ciągła, wybieramy otwarte $U_0\subseteq U_1$ zawierające $x_0$ tak, by $G(x,0)\in V_0$ dla każdego $x\in U_0$. Połóżmy $\varphi(x)=G(x,0)$. Jest to funkcja klasy $C^1$, jako złożenie odwrotności z ciągłymi odwzorowaniami afinicznymi i rzutem. Wówczas $F(x,\varphi(x))=0$, a $\Psi$ jest iniektywna na $U_0\times V_0\subseteq\mathcal O$. Jeśli $F(x,y)=0$, to $\Psi(x,y)=(x,0)=\Psi(x,\varphi(x))$, zatem $y=\varphi(x)$. To daje równoważność oraz jednoznaczność funkcji $\varphi$, ponieważ dla każdego $x$ istnieje dokładnie jeden taki $y\in V_0$. Różniczkując tożsamość $F(x,\varphi(x))=0$ i używając reguły łańcuchowej, otrzymujemy
$D_xF(x,\varphi(x))+D_yF(x,\varphi(x))D\varphi(x)=0$.
Po ewentualnym zmniejszeniu $U_0$, z ciągłości $D_yF$ i $\varphi$ wynika $\|A^{-1}(D_yF(x,\varphi(x))-A)\|<1$; szereg Neumanna daje odwracalność $D_yF(x,\varphi(x)):Y\to Z$. Stąd wzór na $D\varphi$. Wersja wyższej regularności wynika z Wniosku~\ref{cor:inverse-ck15}.
\end{proof}

\begin{przyklad}[Rodzina równań operatorowych]
\label{ex:implicit-operator15}
Niech $E$ będzie rzeczywistą przestrzenią Banacha, $K\in\mathcal L(E)$ i $g\in E$. Na $\R\times E$ definiujemy $F(\lambda,u)=u-\lambda Ku-g$. W punkcie $(0,g)$ mamy $F(0,g)=0$ oraz $D_uF(0,g)=I$, zatem istnieje lokalnie jednoznaczne rozwiązanie $u(\lambda)$ klasy $C^\infty$ równania $(I-\lambda K)u(\lambda)=g$. Z wzoru na pochodną funkcji uwikłanej wynika
$u'(\lambda)=(I-\lambda K)^{-1}Ku(\lambda)$, w szczególności $u'(0)=Kg$. Różniczkując tożsamość $(I-\lambda K)u'(\lambda)=Ku(\lambda)$, otrzymujemy $(I-\lambda K)u''(\lambda)=2Ku'(\lambda)$, więc $u''(0)=2K^2g$. Dla $|\lambda|\,\|K\|<1$ rozwiązanie można również zapisać w postaci zbieżnego szeregu Neumanna $u(\lambda)=\sum_{n\ge0}\lambda^nK^ng$.
\end{przyklad}


\subsection{Różniczkowanie odwracania i eksponenty operatorowej}

\begin{twierdzenie}[Pochodne odwzorowania odwracania]
\label{thm:inverse-derivatives15}
Niech $X$ będzie przestrzenią Banacha nad $\K\in\{\R,\C\}$,
a $GL(X)$ zbiorem odwracalnych elementów algebry
$\mathcal L(X)$. Odwzorowanie $\operatorname{inv}(A)=A^{-1}$
jest klasy $C^\infty$ na $GL(X)$, przy czym w przypadku zespolonym różniczkujemy w sensie zespolonym Frécheta (kolejne pochodne utożsamiamy z operatorami wieloliniowymi nad $\K$). Dla każdego $A\in GL(X)$, $m\ge1$ i $H_1,\ldots,H_m\in\mathcal L(X)$ zachodzi
\[
 D^m\operatorname{inv}(A)[H_1,\ldots,H_m]
 =(-1)^m\sum_{\sigma\in S_m}
 A^{-1}H_{\sigma(1)}A^{-1}\cdots H_{\sigma(m)}A^{-1}.
\]
Tu $S_m$ oznacza zbiór permutacji $\{1,\ldots,m\}$.
\end{twierdzenie}
\begin{proof}
Jeżeli $\|A^{-1}H\|<1$, szereg Neumanna zapewnia $A+H\in GL(X)$, co dowodzi otwartości $GL(X)$ nad obydwoma ciałami, oraz daje
$(A+H)^{-1}=\sum_{k\ge0}(-A^{-1}H)^k A^{-1}$.
Suma wyrazów od $k=2$ ma normę co najwyżej
$\|A^{-1}\|\|A^{-1}H\|^2/(1-\|A^{-1}H\|)$.
Stąd ciągłość odwrotności i wzór
$D\operatorname{inv}(A)[H]=-A^{-1}HA^{-1}$.
Pochodna jest ciągła w normie operatorowej, dzięki
ciągłości mnożenia w algebrze. Indukcyjnie różniczkujemy
każdy czynnik $A^{-1}$ w formule rzędu $m$ regułą
iloczynu. Wstawienie $-A^{-1}H_{m+1}A^{-1}$ w miejsce
jednego z $m+1$ czynników $A^{-1}$ odpowiada wstawieniu
indeksu $m+1$ w jedną z $m+1$ pozycji permutacji.
Każda permutacja $m+1$ indeksów powstaje dokładnie raz.
Otrzymana suma jest ciągła jako odwzorowanie o wartościach
w przestrzeni ograniczonych operatorów wieloliniowych,
bo normę każdego składnika ogranicza iloczyn norm.
To dowodzi indukcyjnie wzoru i gładkości.
\end{proof}

\begin{twierdzenie}[Pochodne eksponenty operatorowej]
\label{thm:exp-derivatives15}
Niech $X$ będzie przestrzenią Banacha nad $\K\in\{\R,\C\}$.
Eksponenta $\exp A=\sum_{n\ge0}A^n/n!$ na $\mathcal L(X)$
jest klasy $C^\infty$ (w sensie zespolonym Frécheta, gdy $\K=\C$). Dla każdego $A\in\mathcal L(X)$, $m\ge1$ i $H_1,\ldots,H_m\in\mathcal L(X)$ zachodzi
\[
 D^m\exp(A)[H_1,\ldots,H_m]
 =\sum_{\sigma\in S_m}\int_{\Delta_m}
 e^{t_0 A} H_{\sigma(1)}e^{t_1 A}\cdots
 H_{\sigma(m)}e^{t_m A}\,\mathrm dt_1\cdots\mathrm dt_m,
\]
gdzie $S_m$ jest zbiorem permutacji $\{1,\ldots,m\}$, $\Delta_m=\{(t_1,\ldots,t_m)\in\R^m:t_j\ge0,\,
\sum_{j=1}^m t_j\le1\}$ i $t_0=1-\sum_{j=1}^m t_j$.
Całkowanie odbywa się względem $m$-wymiarowej miary Lebesgue'a ograniczonej do sympleksu. Można zdefiniować tę całkę przez iterowane całki Riemanna po zmiennych $t_1,\ldots,t_m$ z granicami wynikającymi z nierówności $\sum_jt_j\le1$; w dowodzie wykorzystujemy tę postać i całkowanie jednostajnie zbieżnego szeregu ciągłych funkcji. W szczególności dla $H\in\mathcal L(X)$
$D\exp(A)[H]=\int_0^1 e^{(1-t)A}He^{tA}\,\mathrm dt$.
\end{twierdzenie}
\begin{proof}
Pochodna rzędu $m$ wielomianu $A^n$ jest sumą wszystkich
słów z $n-m$ czynnikami $A$ oraz po jednym czynniku
$H_1,\ldots,H_m$ w dowolnej kolejności. Wynika to
indukcyjnie z reguły iloczynu: w każdym słowie zastępujemy
jeden z pozostałych czynników $A$ nowym kierunkiem.
Dla $n<m$ pochodna jest zerowa, a dla $n\ge m$ liczba słów wynosi $n(n-1)\cdots(n-m+1)=n!/(n-m)!$. Na kuli $\|A\|\le R$, gdzie $R>0$, i dla kierunków $\|H_j\|\le1$ norma każdego słowa nie przekracza $R^{n-m}$. Stąd norma pochodnej $A^n/n!$ jest nie większa niż $R^{n-m}/(n-m)!$.
Zatem dla każdego $m\ge0$ szereg pochodnych zbiega jednostajnie w normie operatorów wieloliniowych na każdej ograniczonej kuli, do ciągłej funkcji $P_m$; dla $m=0$ jest to szereg eksponenty. Jeżeli $P_{m,N}$ oznacza odpowiednią pochodną sumy częściowej, to wzór całkowy na przyrost (w przypadku zespolonym stosowany do krzywej z rzeczywistym parametrem) daje
\[
 P_{m,N}(A+H)-P_{m,N}(A)=\int_0^1 P_{m+1,N}(A+tH)[H,\cdot]\,\mathrm dt .
\]
Dla $m=0$ zapis $[H,\cdot]$ oznacza po prostu ewaluację na $H$. Jednostajna zbieżność na kuli zawierającej odcinek pozwala przejść do granicy. Po odjęciu $P_{m+1}(A)[H,\cdot]$ norma reszty nie przekracza $\|H\|\sup_{0\le t\le1}\|P_{m+1}(A+tH)-P_{m+1}(A)\|=o(\|H\|)$. Ponieważ $P_{m+1}(A)$ jest wieloliniowy nad $\K$, otrzymujemy pochodną Frécheta nad tym samym ciałem. Indukcja po $m$ uzasadnia wszystkie pochodne.
Grupując słowa według liczby $k_j$ czynników $A$
pomiędzy kierunkami, dostajemy
\[
 \sum_{\sigma\in S_m}\sum_{k_0,\ldots,k_m\ge0}
 \frac{A^{k_0}H_{\sigma(1)}A^{k_1}\cdots H_{\sigma(m)}A^{k_m}}
 {(m+k_0+\cdots+k_m)!}.
\]
Dla ustalonej permutacji rozwinięcie eksponent w całce na sympleksie jest jednostajnie bezwzględnie zbieżne w normie: supremum normy składnika o indeksach $(k_0,\ldots,k_m)$ nie przekracza $\prod_j\|H_j\|\prod_{i=0}^m\|A\|^{k_i}/k_i!$, bo $0\le t_i\le1$. Suma tych ograniczeń wynosi $e^{(m+1)\|A\|}\prod_j\|H_j\|$. Uzasadnia to całkowanie szeregu wyraz po wyrazie.
Dla całkowitych $k_j\ge0$ całka jednomianu
$t_0^{k_0}\cdots t_m^{k_m}$ po $\Delta_m$ wynosi
$\prod_{j=0}^m k_j!/(m+\sum_j k_j)!$.
Istotnie, dla $m=1$ wzór wynika z całkowania przez
części całki $\int_0^1 t^p(1-t)^q\,\mathrm dt$,
redukującego $p$ o jeden i zwiększającego $q$ o jeden,
z przypadkiem $p=0$ równym $1/(q+1)$. Dokładniej,
$I(p,q)=pI(p-1,q+1)/(q+1)$ dla $p\ge1$, skąd
$I(p,q)=p!q!/(p+q+1)!$.
W kroku indukcji ustalamy $t_m=u$ i podstawiamy
$t_j=(1-u)v_j$ dla $1\le j<m$. Dla $0\le u<1$
jakobian wynosi $(1-u)^{m-1}$, a $t_0=(1-u)(1-\sum_{j=1}^{m-1}v_j)$.
Przypadek $u=1$ nie wpływa na całkę. Oznaczmy
$P=m-1+\sum_{j=0}^{m-1}k_j$.
Z założenia indukcyjnego całka po $\Delta_{m-1}$,
przed uwzględnieniem czynników zależnych od $u$,
wynosi $\prod_{j=0}^{m-1}k_j!/P!$.
Mnożymy ją przez $\int_0^1u^{k_m}(1-u)^P\,\mathrm du
=k_m!P!/(k_m+P+1)!$ i otrzymujemy
$\prod_{j=0}^{m}k_j!/(m+\sum_{j=0}^m k_j)!$.
Całkowanie rozwinięcia wyraz po wyrazie daje dokładnie
powyższy szereg pochodnej.
\end{proof}

\begin{przyklad}[Nieprzemienność w pochodnej eksponenty]
\label{ex:noncommuting-exp15}
W algebrze macierzy $2\times2$ nad $\R$ lub $\C$ weźmy $A=\operatorname{diag}(0,1)$ i $H=E_{12}$ (jedynka w pozycji $(1,2)$, pozostałe elementy zerowe). Wtedy $e^{(1-t)A}He^{tA}=e^t H$. Zatem
$D\exp(A)[H]=(e-1)H$, podczas gdy $e^AH=H$ i
$He^A=eH$. Jeżeli $AH=HA$, całkowany operator
jest stale równy $e^AH$, więc w tym szczególnym
przypadku $D\exp(A)[H]=e^AH$.
\end{przyklad}

\begin{uwaga}[Dlaczego przestrzenie Banacha i dlaczego pochodna Frécheta?]
\label{rem:limits-inverse15}
Dowód twierdzenia o funkcji odwrotnej wykorzystuje zupełność domkniętej kuli (zasada kontrakcji) oraz jednostajną kontrolę operatora $Df(x)$ w normie. Sama odwracalność pochodnej Gâteaux w punkcie, bez stosownych założeń o ciągłości i pełnej różniczkowalności, nie daje takiego wniosku. W ogólnych przestrzeniach lokalnie wypukłych istnieją odrębne twierdzenia o funkcji odwrotnej, lecz nie wolno przenosić powyższego dowodu na dowolną topologiczną przestrzeń wektorową przez formalne zastąpienie normy seminormami.
\end{uwaga}

\section{Rozmaitości modelowane na przestrzeniach Banacha}
\label{sec:fsg-banach-manifolds}

\subsection{Modele lokalne i zgodność map}

Przestrzeń Banacha jest naturalnym modelem lokalnym przestrzeni
funkcji. W odróżnieniu od przestrzeni euklidesowej może mieć
nieskończony wymiar, a jej domknięte podprzestrzenie nie muszą
mieć domkniętych dopełnień. Dlatego w dalszej teorii istotne
będzie nie tylko istnienie pochodnej, lecz również struktura
jej jądra i obrazu.

Ustalamy rzeczywistą przestrzeń Banacha $E$ oraz
$k\in\{1,2,\ldots,\infty\}$. Klasę $C^k$ odwzorowań między
otwartymi podzbiorami przestrzeni Banacha rozumiemy w sensie
Definicji~\ref{def:ck15}.

\begin{definicja}[Mapa i atlas]
\label{def:fsg-chart}
Niech $M$ będzie przestrzenią topologiczną Hausdorffa.
Mapą o modelu $E$ nazywamy parę $(U,\varphi)$, gdzie
$U\subseteq M$ jest zbiorem otwartym, a
$\varphi:U\to V$ jest homeomorfizmem na otwarty podzbiór
$V\subseteq E$. Dwie mapy $(U,\varphi)$ i $(W,\psi)$
są zgodne klasy $C^k$, jeżeli odwzorowania przejścia
\[
 \psi\circ\varphi^{-1}:\varphi(U\cap W)\longrightarrow
 \psi(U\cap W)
\]
oraz jego odwrotność są klasy $C^k$.
Atlas klasy $C^k$ jest rodziną parami zgodnych map,
których dziedziny pokrywają $M$.
Rozmaitością Banacha klasy $C^k$ nazywamy przestrzeń
$M$ wraz z maksymalnym takim atlasem.
Nie zakładamy w tej definicji drugiego aksjomatu przeliczalności
ani parazwartości; jeżeli będą potrzebne, wskażemy je osobno.
\end{definicja}

\begin{stwierdzenie}[Maksymalne rozszerzenie atlasu]
\label{prop:fsg-maximal-atlas}
Każdy atlas klasy $C^k$ wyznacza dokładnie jeden
maksymalny atlas klasy $C^k$, który go zawiera.
\end{stwierdzenie}
\begin{proof}
Do wyjściowego atlasu $\mathcal A$ dołączamy wszystkie mapy
zgodne z każdą mapą z $\mathcal A$. Aby sprawdzić zgodność
dwóch dołączonych map $(U,\varphi)$, $(W,\psi)$,
ustalmy $p\in U\cap W$. Wybieramy mapę $(O,\theta)\in\mathcal A$
taką, że $p\in O$. Na otwartym otoczeniu $\varphi(U\cap W\cap O)$
punktu $\varphi(p)$ mamy
\[
 \psi\circ\varphi^{-1}
 =(\psi\circ\theta^{-1})\circ(\theta\circ\varphi^{-1}).
\]
Oba czynniki są klasy $C^k$, więc ich złożenie też jest tej
klasy. Analogicznie postępujemy z przejściem odwrotnym.
Ponieważ własność $C^k$ jest lokalna, otrzymujemy zgodność
na całym przecięciu. Maksymalność i jednoznaczność wynikają
z tego, że każda mapa atlasu zawierającego $\mathcal A$
musi być zgodna z każdą mapą z $\mathcal A$.
\end{proof}

\begin{przyklad}[Otwarte podzbiory przestrzeni funkcyjnych]
\label{ex:fsg-positive-functions}
Niech $K$ będzie niepustą zwartą przestrzenią Hausdorffa,
a $E=C(K;\R)$ z normą supremum. Zbiór
\[
 \mathcal P=\{f\in E:f(x)>0\text{ dla każdego }x\in K\}
\]
jest otwarty. Istotnie, dla $f\in\mathcal P$ mamy
$m=\min_{x\in K}f(x)>0$. Jeśli $\|g-f\|_\infty<m/2$,
to $g(x)>m/2$ dla każdego $x\in K$. Mapa tożsamościowa
na $\mathcal P$ wyznacza atlas klasy $C^\infty$.
Jest to rozmaitość Banacha, mimo że jej punkty są funkcjami,
a nie wektorami o skończonej liczbie współrzędnych.
\end{przyklad}

\subsection{Odwzorowania i przestrzenie styczne}

\begin{definicja}[Odwzorowanie klasy $C^r$]
\label{def:fsg-smooth-map}
Niech $M,N$ będą rozmaitościami Banacha klasy $C^k$,
modelowanymi odpowiednio na $E,F$, i niech $1\le r\le k$.
Ciągłe odwzorowanie $f:M\to N$ jest klasy $C^r$, jeżeli
dla każdej pary map $(U,\varphi)$, $(W,\psi)$
odwzorowanie
$\psi\circ f\circ\varphi^{-1}$ jest klasy $C^r$
na otwartym zbiorze $\varphi(U\cap f^{-1}(W))$.
\end{definicja}

Wystarczy sprawdzić ten warunek lokalnie przy jednej parze
map wokół każdego punktu i jego obrazu: dla innych map
lokalny zapis zmienia się przez złożenie z odwzorowaniami
przejścia klasy $C^k$. Reguła łańcuchowa
(Twierdzenie~\ref{thm:chain14}) i jej iteracje zapewniają
zachowanie klasy $C^r$.

\begin{definicja}[Wektor styczny jako klasa krzywych]
\label{def:fsg-tangent}
Niech $p\in M$. Rozważamy krzywe klasy $C^1$
$\gamma:(-\varepsilon,\varepsilon)\to M$ z $\gamma(0)=p$.
Dwie takie krzywe uznajemy za równoważne w punkcie $p$,
jeśli dla pewnej mapy $(U,\varphi)$ wokół $p$
\[
 (\varphi\circ\gamma_1)'(0)
 =(\varphi\circ\gamma_2)'(0).
\]
Przestrzeń styczna $T_pM$ jest zbiorem klas równoważności
tych krzywych.
\end{definicja}

\begin{twierdzenie}[Niezależność od mapy i struktura liniowa]
\label{thm:fsg-tangent-chart}
Relacja z Definicji~\ref{def:fsg-tangent} nie zależy od
wyboru mapy. Dla każdej mapy $(U,\varphi)$ wokół $p$
odwzorowanie
\[
 I_{\varphi,p}:T_pM\to E,\qquad
 I_{\varphi,p}([\gamma])=(\varphi\circ\gamma)'(0)
\]
jest bijekcją. Przeniesienie struktury przestrzeni Banacha
przez tę bijekcję daje strukturę liniową i topologię
niezależne od mapy, choć konkretna norma może zależeć od mapy.
\end{twierdzenie}
\begin{proof}
Dla drugiej mapy $(W,\psi)$ reguła łańcuchowa daje
\[
 (\psi\circ\gamma)'(0)
 =D(\psi\circ\varphi^{-1})(\varphi(p))
   (\varphi\circ\gamma)'(0).
\]
Pochodna przejścia jest ograniczonym izomorfizmem liniowym:
reguła łańcuchowa dla przejścia i jego odwrotności daje
obie tożsamości operatorowe. Równość prędkości w jednej
mapie jest więc równoważna równości w drugiej, co również
dowodzi, że zdefiniowano relację równoważności.
Wstrzykiwanie $I_{\varphi,p}$ wynika z definicji klas.
Dla $v\in E$ wybieramy $\varepsilon>0$ tak, by
$\varphi(p)+tv\in\varphi(U)$ dla $|t|<\varepsilon$.
Krzywa $\gamma(t)=\varphi^{-1}(\varphi(p)+tv)$ ma
współrzędną prędkość $v$, co daje surjekcję.
Zmiana mapy odpowiada ograniczonemu izomorfizmowi liniowemu,
zatem przeniesione działania liniowe są identyczne, a normy
równoważne. Zupełność wynika z zupełności $E$.
\end{proof}

\begin{twierdzenie}[Różniczka odwzorowania rozmaitości]
\label{thm:fsg-tangent-map}
Dla odwzorowania $f:M\to N$ klasy $C^1$ różniczka
\[
 T_pf:T_pM\to T_{f(p)}N,\qquad [\gamma]\mapsto[f\circ\gamma]
\]
jest dobrze określonym ograniczonym operatorem liniowym.
W mapach $\varphi,\psi$ wokół $p,f(p)$ zachodzi
\[
 I_{\psi,f(p)}\circ T_pf\circ I_{\varphi,p}^{-1}
 =D(\psi\circ f\circ\varphi^{-1})(\varphi(p)).
\]
Jeżeli $g:N\to P$ jest klasy $C^1$, to
$T_p(g\circ f)=T_{f(p)}g\circ T_pf$.
\end{twierdzenie}
\begin{proof}
Współrzędna prędkość $f\circ\gamma$ jest, z reguły
łańcuchowej, obrazem prędkości $\varphi\circ\gamma$
przez operator po prawej stronie równości. Zależy więc
wyłącznie od klasy $[\gamma]$. Wzór współrzędny
daje liniowość i ograniczoność. Ostatnia równość wynika
z $(g\circ f)\circ\gamma=g\circ(f\circ\gamma)$.
\end{proof}

\begin{przyklad}[Różniczka funkcjonału całkowego]
\label{ex:fsg-integral-tangent}
Na rozmaitości $\mathcal P\subset C(K;\R)$ z
Przykładu~\ref{ex:fsg-positive-functions} niech $\mu$
będzie skończoną dodatnią miarą borelowską na $K$.
Funkcjonał $J(f)=\int_K f(x)^2\,\mathrm d\mu(x)$
ma różniczkę
\[
 T_fJ(h)=2\int_K f(x)h(x)\,\mathrm d\mu(x),
 \qquad h\in T_f\mathcal P\simeq C(K;\R).
\]
Istotnie, reszta po odjęciu części liniowej jest równa
$\int_K h^2\,\mathrm d\mu$, a jej moduł nie przekracza
$\mu(K)\|h\|_\infty^2$. Operator liniowy ma normę
nie większą niż $2\mu(K)\|f\|_\infty$.
Ponadto zmiana tego operatora przy zmianie $f$ na $g$
ma normę nie większą niż $2\mu(K)\|f-g\|_\infty$,
więc $J$ jest klasy $C^1$.
\end{przyklad}

% Wersja źródłowa: audyt II, jednostka 16

\section{Funkcjonały wypukłe i podstawy rachunku wariacyjnego}
\label{sec:wariacyjny16}

\providecommand{\atlasref}[2]{#2}


\subsection{Od pochodnej Frécheta do problemu minimum}

\noindent\emph{Powiązanie tematyczne w Bibliotece matematycznej:} \atlasref{AM1-25::ma1:def:convex-function}{AM1-25: wypukłość funkcji jednej zmiennej}.


W poprzednich rozdziałach badaliśmy lokalną strukturę funkcji
różniczkowalnych na przestrzeniach Banacha. Rachunek wariacyjny pyta
o istnienie i charakter minimizatorów funkcjonału $J:E\to\R$,
często na nieskończenie wymiarowej przestrzeni funkcji.
W tym rozdziale $H$ oznacza rzeczywistą przestrzeń Hilberta,
a $E$ rzeczywistą przestrzeń Banacha. Pisząc $\langle\cdot,\cdot\rangle$
w $H$, przyjmujemy liniowość w pierwszym argumencie.

\begin{definicja}[Wypukłość funkcjonału]
\label{def:convex-fun16}
Niech $E$ będzie rzeczywistą przestrzenią liniową,
a $C\subseteq E$ zbiorem wypukłym.
Funkcjonał $J:C\to(-\infty,+\infty]$ nazywamy
\emph{wypukłym}, jeżeli
$J((1-t)x+ty)\le(1-t)J(x)+tJ(y)$ dla wszystkich $x,y\in C$
i $0<t<1$. Dla $t=0,1$ przyjmujemy równość, bez
mnożenia zera przez nieskończoność.
Funkcjonał nazywamy \emph{właściwym}, gdy jego dziedzina
efektywna $\operatorname{dom}J=\{x\in C:J(x)<+\infty\}$
jest niepusta. Z wypukłości $J$ wynika wypukłość
$\operatorname{dom}J$, ponieważ kombinacja wypukła
dwóch skończonych wartości jest skończona.
Funkcjonał wypukły jest \emph{ściśle wypukły}, jeśli nierówność
jest ostra dla różnych $x,y\in\operatorname{dom}J$
i $0<t<1$. W przypadku funkcjonału rzeczywistego
$\operatorname{dom}J=C$.
\end{definicja}

\begin{twierdzenie}[Nierówność styczna]
\label{thm:convex-derivative16}
Niech $E$ będzie rzeczywistą przestrzenią Banacha,
$U\subseteq E$ niepustym otwartym zbiorem wypukłym, a
$J:U\to\R$ funkcjonałem różniczkowalnym w sensie
Frécheta w każdym punkcie $U$.
Następujące warunki są równoważne:
\begin{enumerate}
\item $J$ jest wypukły na $U$;
\item $J(y)\ge J(x)+DJ(x)(y-x)$ dla wszystkich $x,y\in U$.
\end{enumerate}
Jeżeli $DJ(x_0)=0$, gdzie $x_0\in U$, a $J$ jest wypukły,
to $x_0$ jest globalnym minimizatorem $J$ na $U$.
Odwrotnie, jeśli $x_0$ jest lokalnym minimizatorem
różniczkowalnego $J$ i $x_0\in U$, to $DJ(x_0)=0$.
\end{twierdzenie}
\begin{proof}
Dla ustalonych $x,y\in U$ niech
$\phi(t)=J(x+t(y-x))$, $0\le t\le1$.
Jeżeli $J$ jest wypukły, to dla $0<t\le1$
z definicji wypukłości
$[\phi(t)-\phi(0)]/t\le\phi(1)-\phi(0)$.
Z reguły łańcuchowej (Twierdzenie~\ref{thm:chain14})
pochodna $\phi$ w zerze wynosi $DJ(x)(y-x)$
(otwartość $U$ pozwala przedłużyć parametr nieco poza zero).
Przechodząc z $t\downarrow0$, otrzymujemy
$DJ(x)(y-x)\le J(y)-J(x)$.

Odwrotnie, ustalmy $x,y\in U$, $0<t<1$ i
$z=(1-t)x+ty\in U$.
Nierówność z punktu 2 daje
$J(x)\ge J(z)+DJ(z)(x-z)$ oraz
$J(y)\ge J(z)+DJ(z)(y-z)$.
Mnożąc je odpowiednio przez $1-t$ i $t$ oraz dodając,
otrzymujemy
\[
 (1-t)J(x)+tJ(y)\ge J(z)+DJ(z)((1-t)x+ty-z)=J(z).
\]
Dla $t=0,1$ zachodzi równość, więc $J$ jest wypukły.
Jeżeli $DJ(x_0)=0$, nierówność styczna zapewnia
$J(y)\ge J(x_0)$. Jeśli zaś $x_0$ jest lokalnym minimum,
dla każdego $h\in E$ i dostatecznie małych $|t|$
punkt $x_0+th$ należy do $U$ i do otoczenia lokalnego
minimum. Funkcja $t\mapsto J(x_0+th)$ ma zatem
minimum w zerze, a jej pochodna w zerze, równa
$DJ(x_0)h$, musi znikać. Ponieważ dotyczy to
każdego $h\in E$, mamy $DJ(x_0)=0$.
\end{proof}

\begin{definicja}[Silna wypukłość]
\label{def:strong-convex16}
Niech $H$ będzie rzeczywistą przestrzenią Hilberta
i $\mu>0$. Funkcjonał $J:H\to\R$ jest
$\mu$-silnie wypukły, jeśli $x\mapsto J(x)-\frac\mu2\|x\|^2$
jest wypukły.
\end{definicja}

\begin{stwierdzenie}[Koercywność z silnej wypukłości]
\label{prop:strong-coercive16}
Niech $H$ będzie rzeczywistą przestrzenią Hilberta,
$\mu>0$, a $J:H\to\R$ funkcjonałem $\mu$-silnie wypukłym
i różniczkowalnym w sensie Frécheta na $H$. Wtedy dla $x,y\in H$
zachodzi
\[
J(y)\ge J(x)+DJ(x)(y-x)+\frac\mu2\|y-x\|^2.
\]
W szczególności $J$ jest koercywny, tj.
$J(y)\to+\infty$, gdy $\|y\|\to\infty$,
i ma co najwyżej jeden minimizator.
\end{stwierdzenie}
\begin{proof}
Stosujemy Twierdzenie~\ref{thm:convex-derivative16}
do funkcjonału $G(z)=J(z)-\mu\|z\|^2/2$.
Z rozwinięcia
$\|x+h\|^2=\|x\|^2+2\langle x,h\rangle+\|h\|^2$
wynika $DG(x)h=DJ(x)h-\mu\langle x,h\rangle$.
Nierówność styczna dla $G$ daje
\[
 J(y)-\tfrac\mu2\|y\|^2
 \ge J(x)-\tfrac\mu2\|x\|^2
       +DJ(x)(y-x)-\mu\langle x,y-x\rangle .
\]
Po przeniesieniu składników kwadratowych na prawą stronę
i wykorzystaniu tożsamości
$\|y\|^2-\|x\|^2-2\langle x,y-x\rangle=\|y-x\|^2$
otrzymujemy żądaną nierówność.
Dla $x=0$ mamy
$J(y)\ge J(0)-\|DJ(0)\|\|y\|+\mu\|y\|^2/2$,
co zapewnia koercywność. Jeżeli $x\ne y$ są minimizatorami,
to wypukłość $G$ przy $t=1/2$ i tożsamość
$\tfrac12(\|x\|^2+\|y\|^2)-\|(x+y)/2\|^2
=\|x-y\|^2/4$ dają
$J((x+y)/2)\le [J(x)+J(y)]/2-\mu\|x-y\|^2/8$,
sprzecznie z minimalnością.
\end{proof}

\begin{przyklad}[Jednostajne dolne oszacowanie hesjanu i silna wypukłość]
\label{ex:variational-hessian16}
Niech $H$ będzie rzeczywistą przestrzenią Hilberta
i $J\in C^2(H,\R)$. Jeżeli istnieje $\mu>0$ takie, że
$D^2J(x)[h,h]\ge\mu\|h\|^2$ dla wszystkich $x,h\in H$,
to $J$ jest $\mu$-silnie wypukły.
Istotnie, dla $x,y\in H$ i $h=y-x$ funkcja
$\phi(t)=J(x+th)-\frac\mu2\|x+th\|^2$
ma drugą pochodną
$\phi''(t)=D^2J(x+th)[h,h]-\mu\|h\|^2\ge0$.
Zatem $\phi'$ jest niemalejąca. Dla $0<t<1$
z twierdzenia o wartości średniej na $[0,t]$ i $[t,1]$
wynika
$(\phi(t)-\phi(0))/t\le(\phi(1)-\phi(t))/(1-t)$.
Po pomnożeniu przez $t(1-t)$ i uporządkowaniu
otrzymujemy $\phi(t)\le(1-t)\phi(0)+t\phi(1)$,
czyli wypukłość funkcjonału $J-\mu\|\cdot\|^2/2$.
Sama dodatnia określoność $D^2J(0)$, nawet wraz
z $DJ(0)=0$, nie gwarantuje lokalnego minimum:
\hyperref[ex:positive-hessian15]{kontrprzykład w rozdziale 15}
podaje dokładnie taki funkcjonał na $\ell^2$.
\end{przyklad}

\subsection{Zbieżność słaba i metoda bezpośrednia}

\begin{definicja}[Półciągłość z dołu i koercywność]
\label{def:weak-lsc16}
Niech $H$ będzie rzeczywistą przestrzenią Hilberta,
$C\subseteq H$ niepustym zbiorem i $J:C\to(-\infty,+\infty]$.
Funkcjonał $J$ jest \emph{słabo półciągły z dołu}, jeśli dla każdej sieci
$(x_\alpha)$ w $C$ zbieżnej do $x\in C$ w topologii
słabej $\sigma(H,H')$ zachodzi
$J(x)\le\liminf_\alpha J(x_\alpha)$, gdzie
$\liminf_\alpha J(x_\alpha)=
\sup_\alpha\inf_{\beta\ge\alpha}J(x_\beta)$.
Równoważnie, każdy zbiór podpoziomowy
$\{x\in C:J(x)\le a\}$, $a\in\R$, jest domknięty
w topologii słabej indukowanej na $C$.
Istotnie, warunek dla sieci zapewnia, że granica sieci
z podpoziomu nadal należy do tego podpoziomu.
Odwrotnie, gdyby $\liminf_\alpha J(x_\alpha)<J(x)$,
można byłoby wybrać rzeczywiste $a$ ściśle pomiędzy
tymi wartościami. Otwarte dopełnienie podpoziomu
zawiera $x$, więc od pewnego indeksu $J(x_\alpha)>a$,
co wymusza $\liminf_\alpha J(x_\alpha)\ge a$,
sprzecznie z wyborem $a$.
Analogicznie definiuje się półciągłość z dołu
w topologii normowej.
Funkcjonał jest \emph{koercywny na $C$}, jeśli dla
każdego $a\in\R$ istnieje $R>0$ takie, że
$J(x)>a$ dla każdego $x\in C$ o $\|x\|>R$.
Równoważnie, wszystkie podpoziomy są ograniczone.
\end{definicja}

\begin{lemat}[Podpoziomowe zbiory wypukłe]
\label{lem:convex-weaklsc16}
Niech $H$ będzie rzeczywistą przestrzenią Hilberta,
a $J:H\to(-\infty,+\infty]$ funkcjonałem wypukłym
i półciągłym z dołu w topologii normowej.
Wtedy $J$ jest słabo półciągły z dołu.
\end{lemat}
\begin{proof}
Dla każdego $a\in\R$ zbiór $C_a=\{x\in H:J(x)\le a\}$
jest normowo domknięty przez półciągłość z dołu
i wypukły przez wypukłość $J$. Z Wniosku~\ref{cor:weak-convex-closure}
o równości domknięć zbioru wypukłego wynika,
że każdy normowo domknięty zbiór wypukły jest
domknięty także w topologii słabej. Zatem wszystkie $C_a$
są słabo domknięte, co daje tezę.
\end{proof}

\begin{lemat}[Słaba zbieżność podciągu w przestrzeni Hilberta]
\label{lem:hilbert-weakseq16}
Niech $H$ będzie rzeczywistą przestrzenią Hilberta,
a $(x_n)$ ciągiem ograniczonym w $H$.
Istnieje podciąg $(x_{n_k})$ i $x\in H$ takie,
że $x_{n_k}\to x$ słabo w $H$.
Nie zakładamy ośrodkowości całej przestrzeni $H$.
\end{lemat}
\begin{proof}
Niech $M=\sup_n\|x_n\|<\infty$ i
$H_0=\overline{\operatorname{lin}\{x_n:n\in\N\}}$.
Przestrzeń $H_0$ jest ośrodkowa, ponieważ skończone
kombinacje liniowe $x_n$ o współczynnikach wymiernych
tworzą w niej przeliczalny zbiór gęsty. Jest również
przestrzenią Hilberta jako domknięta podprzestrzeń $H$.
Jeśli $H_0=\{0\}$, cały ciąg jest zerowy i teza zachodzi.
W przeciwnym razie wybieramy bazę ortonormalną $(e_j)$
skończoną lub przeliczalnie nieskończoną.
Dla każdego $j$ ciąg $\langle x_n,e_j\rangle$ jest ograniczony.
Z twierdzenia Bolzana--Weierstrassa wybieramy podciąg,
na którym pierwsza współrzędna ma granicę. Z niego
wybieramy podciąg dla drugiej współrzędnej i tak dalej.
W przypadku skończonej bazy ostatni podciąg ma
wszystkie wymagane granice. W przypadku nieskończonej
bazy wybieramy $k$-ty wyraz $k$-tego podciągu.
Indeksy tak wybranego ciągu rosną, a jego ogon
od $k=j$ jest podciągiem $j$-tego podciągu, więc
$j$-ta współrzędna jest zbieżna dla każdego $j$.

Na przestrzeni skończonych kombinacji wektorów $e_j$
definiujemy $L(v)=\lim_k\langle x_{n_k},v\rangle$.
Istnienie granicy wynika z jej istnienia na każdym $e_j$,
a $|L(v)|\le M\|v\|$. Funkcjonał $L$ przedłuża się
jednoznacznie na $H_0$ jako ciągły funkcjonał liniowy.
Z twierdzenia Riesza istnieje $x\in H_0$ takie,
że $L(v)=\langle x,v\rangle$ dla $v\in H_0$.
Ponieważ $\|x\|=\|L\|\le M$, dla dowolnego $v\in H_0$
i skończonej kombinacji $w$ wektorów $e_j$ zachodzi
\[
 |\langle x_{n_k}-x,v\rangle|
 \le 2M\|v-w\|+|\langle x_{n_k}-x,w\rangle|.
\]
Dla $\varepsilon>0$ wybieramy $w$ tak, aby pierwszy
składnik był mniejszy od $\varepsilon/2$, a następnie
$k$ dostatecznie duże, aby drugi składnik był mniejszy
od $\varepsilon/2$. Gdy $M=0$, ciąg jest zerowy.
Otrzymujemy zbieżność $\langle x_{n_k},v\rangle
\to\langle x,v\rangle$ dla każdego $v\in H_0$.
Dla $v\in H_0^\perp$ obie strony są zerowe.
Dla dowolnego $v\in H$ stosujemy rozkład
$H=H_0\oplus H_0^\perp$ z Twierdzenia~\ref{thm:orth-decomposition13}.
Z twierdzenia Riesza każdy ciągły funkcjonał liniowy
na $H$ jest dany przez iloczyn skalarny z takim $v$.
Otrzymana zbieżność jest więc zbieżnością słabą w $H$.
\end{proof}

\begin{twierdzenie}[Metoda bezpośrednia w przestrzeni Hilberta]
\label{thm:direct-method16}
Niech $H$ będzie rzeczywistą przestrzenią Hilberta,
$C\subseteq H$ niepustym, normowo domkniętym i wypukłym
zbiorem, a $J:C\to(-\infty,+\infty]$ funkcjonałem,
który jest właściwy (tzn. przyjmuje przynajmniej jedną
wartość skończoną), ograniczony z dołu, koercywny na $C$
i półciągły z dołu dla topologii słabej indukowanej na $C$.
Wtedy istnieje $u\in C$ takie, że
$J(u)=\inf_{v\in C}J(v)$.
Jeżeli dodatkowo $J$ jest ściśle wypukły,
minimizator jest jedyny.
\end{twierdzenie}
\begin{proof}
Niech $m=\inf_C J$, który jest skończony na mocy
właściwości i ograniczoności z dołu. Dla każdego $n\ge1$ z definicji infimum wybieramy
$u_n\in C$ takie, że $m\le J(u_n)<m+1/n$.
Wszystkie $u_n$ należą do podpoziomu $\{v\in C:J(v)\le m+1\}$,
który jest ograniczony z koercywności. Zatem
$(u_n)$ jest ograniczonym ciągiem w $H$.
Na mocy Lematu~\ref{lem:hilbert-weakseq16}
pewien podciąg $u_{n_k}\to u$ słabo w $H$.
Ponieważ $C$ jest normowo domknięty i wypukły,
jest również słabo domknięty na mocy
Wniosku~\ref{cor:weak-convex-closure}, więc $u\in C$.
Słaba półciągłość z dołu daje
$J(u)\le\liminf_k J(u_{n_k})=m$; z definicji infimum
zachodzi też $J(u)\ge m$. Jeśli dodatkowo $J$ jest
ściśle wypukły i $u\ne v$ są minimizatorami, to
$J(u)=J(v)=m$ jest skończone, a $(u+v)/2\in C$.
Zatem $J((u+v)/2)<m$, co przeczy definicji $m$.
\end{proof}

\begin{uwaga}[Czego nie daje sama koercywność?]
\label{rem:coercive16}
W nieskończenie wymiarowej przestrzeni Hilberta
domknięta kula jednostkowa nie jest zwarta w normie:
zawiera ciąg ortonormalny, którego różne wyrazy
są odległe o $\sqrt2$, więc żaden jego podciąg
nie jest zbieżny w normie.
Sama ciągłość normowa i koercywność funkcjonału
nie dają zatem zwartości podpoziomu potrzebnej
do bezpośredniego zastosowania twierdzenia Weierstrassa. W twierdzeniu wykorzystaliśmy
słabą zbieżność podciągów i słabą półciągłość z dołu.
W ogólnej przestrzeni Banacha naturalną rolę przestrzeni
Hilberta przejmują przestrzenie refleksywne.
W rozdziale 26 udowodnimy ogólną wersję twierdzenia
o słabo zbieżnym podciągu ciągu ograniczonego.
Powyższy elementarny dowód w przestrzeni Hilberta
nie wymaga tej późniejszej teorii.
\end{uwaga}

\subsection{Subgradienty i odwzorowanie proksymalne}

\begin{definicja}[Subgradient i subróżniczka]
\label{def:subgrad16}
Niech $H$ będzie rzeczywistą przestrzenią Hilberta,
$J:H\to(-\infty,+\infty]$ właściwym funkcjonałem wypukłym,
a $u\in H$ punktem takim,
że $J(u)<+\infty$. Wektor $p\in H$ nazywamy
\emph{subgradientem} $J$ w $u$, jeśli
$J(v)\ge J(u)+\langle p,v-u\rangle$ dla każdego $v\in H$.
Zbiór wszystkich subgradientów oznaczamy $\partial J(u)$
i nazywamy \emph{subróżniczką} $J$ w $u$.
Dla $J(u)=+\infty$ przyjmujemy $\partial J(u)=\varnothing$.
\end{definicja}

\begin{stwierdzenie}[Warunek optymalności i monotoniczność]
\label{prop:subgrad16}
Niech $H$ będzie rzeczywistą przestrzenią Hilberta,
a $J:H\to(-\infty,+\infty]$ właściwym funkcjonałem wypukłym.
Dla każdego $u\in H$ punkt $u$ minimalizuje $J$ wtedy
i tylko wtedy, gdy $0\in\partial J(u)$. Jeśli $u,v\in H$,
$p\in\partial J(u)$ oraz $q\in\partial J(v)$, to
$\langle p-q,u-v\rangle\ge0$.
Jeżeli dla $u\in H$ funkcjonał $J$ jest skończony
na otoczeniu $u$ i różniczkowalny w sensie Frécheta w $u$,
to $\partial J(u)=\{\nabla J(u)\}$, gdzie $\nabla J(u)\in H$
jest jedynym wektorem spełniającym
$DJ(u)h=\langle\nabla J(u),h\rangle$ dla każdego $h\in H$.
\end{stwierdzenie}
\begin{proof}
Jeżeli $u$ minimalizuje $J$, to $J(u)$ jest skończone,
gdyż $J$ jest właściwy. Nierówność $J(v)\ge J(u)$
dla każdego $v\in H$ oznacza dokładnie, że $0\in\partial J(u)$.
Odwrotnie, przynależność $0\in\partial J(u)$ daje tę
nierówność, a więc globalną minimalność.
Z nierówności subgradientowych
$J(v)\ge J(u)+\langle p,v-u\rangle$ i
$J(u)\ge J(v)+\langle q,u-v\rangle$,
po dodaniu i skróceniu skończonych wartości $J(u),J(v)$
otrzymujemy $0\ge\langle p,v-u\rangle+\langle q,u-v\rangle
=-\langle p-q,u-v\rangle$, czyli monotoniczność.
Jeśli $J$ jest różniczkowalny, to dla $p\in\partial J(u)$
dla $h\in H$ podstawiamy $v=u+th$ przy dostatecznie małym
$|t|>0$. Po podzieleniu nierówności przez $t>0$ mamy
$(J(u+th)-J(u))/t\ge\langle p,h\rangle$, zaś dla $t<0$
nierówność ma przeciwny zwrot. Oba ilorazy różnicowe
dążą do $DJ(u)h$, więc $\langle p,h\rangle=DJ(u)h$.
Ponieważ $h\in H$ było dowolne, z jednoznaczności
reprezentacji Riesza wynika $p=\nabla J(u)$.
Odwrotnie, dla $v$ z dziedziny efektywnej $J$ i $0<t<1$
wypukłość daje
$(J(u+t(v-u))-J(u))/t\le J(v)-J(u)$.
Przechodząc z $t$ do zera, otrzymujemy
$DJ(u)(v-u)\le J(v)-J(u)$.
Dla $J(v)=+\infty$ nierówność zachodzi automatycznie.
Reprezentacja Riesza daje zatem $\nabla J(u)\in\partial J(u)$.
\end{proof}

\begin{przyklad}[Subróżniczka normy w Hilbercie]
\label{ex:norm-subgrad16}
Niech $H$ będzie rzeczywistą przestrzenią Hilberta
i $J:H\to\R$, $J(x)=\|x\|$. Wówczas
$\partial J(0)=\{p:\|p\|\le1\}$ i
$\partial J(x)=\{x/\|x\|\}$ dla $x\ne0$.
Jeżeli $\|p\|\le1$, to $\langle p,v\rangle\le\|v\|$
dla każdego $v\in H$, zatem $p\in\partial J(0)$.
Odwrotnie, dla $p\in\partial J(0)$, podstawiając $v=p$,
otrzymujemy $\|p\|^2\le\|p\|$, czyli $\|p\|\le1$.
Dla $x\ne0$ reguła łańcuchowa zastosowana do
$\|x\|=\sqrt{\langle x,x\rangle}$ daje
$DJ(x)h=\langle x,h\rangle/\|x\|$.
Stwierdzenie~\ref{prop:subgrad16} daje drugi wzór.
W ogólnej przestrzeni Banacha analogiczny subgradient normy
jest funkcjonałem z przestrzeni dualnej, a jego istnienie
dla $x\ne0$ zapewnia twierdzenie Hahna--Banacha.
\end{przyklad}

\begin{twierdzenie}[Odwzorowanie proksymalne]
\label{thm:prox16}
Niech $H$ będzie rzeczywistą przestrzenią Hilberta,
a $J:H\to(-\infty,+\infty]$ właściwym, wypukłym,
normowo półciągłym z dołu i ograniczonym z dołu funkcjonałem,
a $\lambda>0$. Dla każdego $y\in H$ istnieje dokładnie
jeden punkt $u=\operatorname{prox}_{\lambda J}(y)$
minimalizujący funkcjonał
$v\mapsto J(v)+\|v-y\|^2/(2\lambda)$.
Dla każdego $u,y\in H$ zachodzi równoważność
$u=\operatorname{prox}_{\lambda J}(y)$ wtedy i tylko wtedy,
gdy $y-u\in\lambda\partial J(u)$.
Ponadto odwzorowanie
$\operatorname{prox}_{\lambda J}:H\to H$ jest
nieekspansywne:
$\|\operatorname{prox}_{\lambda J}(y)-
\operatorname{prox}_{\lambda J}(z)\|\le\|y-z\|$.
\end{twierdzenie}
\begin{proof}
Funkcjonał $\Phi_y(v)=J(v)+\|v-y\|^2/(2\lambda)$
jest właściwy, wypukły i normowo półciągły z dołu,
a więc słabo półciągły z dołu na mocy
Lematu~\ref{lem:convex-weaklsc16}.
Jeżeli $J(v)\ge m$ dla każdego $v\in H$, to
$\Phi_y(v)\ge m+\|v-y\|^2/(2\lambda)\ge m$.
Nierówność $\|v-y\|\ge\|v\|-\|y\|$ zapewnia koercywność.
Dla różnych $v,w$ z dziedziny efektywnej i $0<t<1$,
z wypukłości $J$ i tożsamości kwadratowej wynika
$\Phi_y((1-t)v+tw)\le(1-t)\Phi_y(v)+t\Phi_y(w)
-t(1-t)\|v-w\|^2/(2\lambda)$.
Zatem $\Phi_y$ jest ściśle wypukły.
Twierdzenie~\ref{thm:direct-method16} dla $C=H$
daje jednoznaczny minimizator $u$.

Jeśli $u$ minimalizuje $\Phi_y$, to dla $v$ z
$J(v)<+\infty$ i $0<t<1$, z wypukłości $J$
oraz minimalności mamy
\[
0\le\Phi_y(u+t(v-u))-\Phi_y(u)
\le t[J(v)-J(u)]+\frac t\lambda\langle u-y,v-u\rangle
+\frac{t^2}{2\lambda}\|v-u\|^2.
\]
Po podzieleniu przez $t$ i przejściu z $t$ do zera
otrzymujemy $J(v)\ge J(u)+\langle(y-u)/\lambda,v-u\rangle$.
Dla $J(v)=+\infty$ nierówność subgradientowa jest
automatyczna, a więc $(y-u)/\lambda\in\partial J(u)$.
Odwrotnie, jeśli $(y-u)/\lambda\in\partial J(u)$,
to $J(u)<+\infty$. Dla $v$ z dziedziny efektywnej
rozwijamy $\|v-y\|^2=\|u-y\|^2+
2\langle u-y,v-u\rangle+\|v-u\|^2$ i dostajemy
\[
\Phi_y(v)-\Phi_y(u)
\ge\langle (y-u)/\lambda,v-u\rangle+
\lambda^{-1}\langle u-y,v-u\rangle+
\frac{\|v-u\|^2}{2\lambda}
=\frac{\|v-u\|^2}{2\lambda}.
\]
Poza dziedziną efektywną $\Phi_y(v)=+\infty$, więc
także tam $\Phi_y(v)\ge\Phi_y(u)$. To zapewnia minimalność.

Niech $u=\operatorname{prox}_{\lambda J}(y)$ oraz
$v=\operatorname{prox}_{\lambda J}(z)$.
Monotoniczność subróżniczki daje
$\langle (y-u)-(z-v),u-v\rangle\ge0$.
Stąd $\|u-v\|^2\le\langle y-z,u-v\rangle
\le\|y-z\|\,\|u-v\|$,
co po podzieleniu przez $\|u-v\|$ (gdy $u\ne v$)
daje tezę. Gdy $u=v$, nierówność jest oczywista.
\end{proof}

\subsection{Przykład: słabe zagadnienie brzegowe}

Niech $\lambda$ oznacza miarę Lebesgue'a na $[0,1]$.
Wprowadzamy przestrzeń $H_0^1(0,1)$ przez opis elementarny:
\[
H_0^1(0,1)=\left\{u:[0,1]\to\R:\ 
\begin{array}{l}
\exists v\in L^2(0,1;\R),\ \int_0^1 v(t)\,\mathrm dt=0,\\
u(x)=\int_0^x v(t)\,\mathrm dt\quad\text{dla każdego }x\in[0,1]
\end{array}\right\}.
\]
Z nierówności Cauchy'ego--Schwarza mamy
$\|v\|_{L^1}\le\|v\|_{L^2}$, zatem całka określająca $u$
jest niezależna od reprezentanta klasy $v$ w $L^2$.
Z twierdzenia o całce nieoznaczonej Lebesgue'a
funkcje $u$ są absolutnie ciągłe i $u'=v$ prawie wszędzie,
a ze wzoru całkowego i warunku zerowej całki wynikają
warunki brzegowe $u(0)=u(1)=0$. Przyjmujemy
$\langle u,w\rangle_H=\int_0^1u'(t)w'(t)\,\mathrm dt$.
Odwzorowanie $u\mapsto u'$ jest liniową izometrią
$H_0^1(0,1)$ na domkniętą podprzestrzeń
$\{v\in L^2(0,1):\int_0^1v(t)\,\mathrm dt=0\}$,
przy czym podprzestrzeń ta jest jądrem ciągłego funkcjonału
$v\mapsto\int_0^1v(t)\,\mathrm dt$ (jego norma jest co najwyżej $1$
na mocy nierówności Cauchy'ego--Schwarza).
Surjektywność izometrii wynika z definicji $H_0^1(0,1)$,
a iniektywność z warunku $u(0)=0$ i wzoru całkowego.
Jako domknięta podprzestrzeń przestrzeni Hilberta $L^2(0,1)$
jej obraz jest zupełny, więc $H_0^1(0,1)$ jest przestrzenią Hilberta.

\begin{lemat}[Nierówność Poincarégo na odcinku]
\label{lem:poincare16}
Dla każdego $u\in H_0^1(0,1)$ zachodzi
$\|u\|_{L^2}\le2^{-1/2}\|u'\|_{L^2}$.
\end{lemat}
\begin{proof}
Ponieważ $u(x)=\int_0^xu'(t)\,\mathrm dt$,
nierówność Cauchy'ego--Schwarza daje
$|u(x)|^2\le x\int_0^x|u'(t)|^2\,\mathrm dt
\le x\|u'\|_{L^2}^2$.
Całkując po $x\in[0,1]$, otrzymujemy
$\|u\|_{L^2}^2\le\frac12\|u'\|_{L^2}^2$.
\end{proof}

\begin{twierdzenie}[Zasada minimalizacji dla problemu Dirichleta]
\label{thm:dirichlet16}
Niech $f\in L^2(0,1;\R)$. Na $H=H_0^1(0,1)$
określmy funkcjonał
\[
J(u)=\frac12\int_0^1|u'(t)|^2\,\mathrm dt
-\int_0^1 f(t)u(t)\,\mathrm dt.
\]
Istnieje dokładnie jeden minimizator $u\in H$ funkcjonału $J$.
Jest on dokładnie jedynym elementem $H$ spełniającym
równanie słabe
\[
\int_0^1u'(t)v'(t)\,\mathrm dt
=\int_0^1f(t)v(t)\,\mathrm dt\qquad(v\in H).
\]
Dla $f(t)=1$ rozwiązaniem jest $u(t)=t(1-t)/2$.
\end{twierdzenie}
\begin{proof}
Niech $\ell(u)=\int_0^1 f(t)u(t)\,\mathrm dt$.
Z Lematu~\ref{lem:poincare16} i nierówności Cauchy'ego--Schwarza
mamy $|\ell(u)|\le2^{-1/2}\|f\|_{L^2}\|u\|_H$,
więc $\ell$ jest ciągłym funkcjonałem liniowym na $H$.
Ponieważ $J(u)=\|u\|_H^2/2-\ell(u)$, jest on ciągły.
Dla $u,w\in H$ i $0<t<1$ liniowość $\ell$ oraz
rozwinięcie kwadratu normy dają
\[
(1-t)J(u)+tJ(w)-J((1-t)u+tw)
=\tfrac12t(1-t)\|u-w\|_H^2.
\]
Prawa strona jest dodatnia dla $u\ne w$, więc $J$
jest ściśle wypukły. Z oszacowania $\ell$ wynika
$J(u)\ge\frac12\|u\|_H^2-
2^{-1/2}\|f\|_{L^2}\|u\|_H$, więc jest koercywny
i ograniczony z dołu (dla $a=2^{-1/2}\|f\|_{L^2}$
i $r=\|u\|_H$ mamy $r^2/2-ar=(r-a)^2/2-a^2/2\ge-a^2/2$).
Z Lematu~\ref{lem:convex-weaklsc16}, ponieważ $J$ jest
wypukły i normowo ciągły, jest słabo półciągły z dołu.
Twierdzenie~\ref{thm:direct-method16} daje minimizator,
a ścisła wypukłość zapewnia jego jednoznaczność.

Dla $h\in H$ rozwinięcie kwadratu daje
\[
J(u+h)-J(u)=\langle u,h\rangle_H-\int_0^1 f(t)h(t)\,\mathrm dt
+\tfrac12\|h\|_H^2.
\]
Część liniowa jest ograniczona przez
$(\|u\|_H+2^{-1/2}\|f\|_{L^2})\|h\|_H$,
a reszta podzielona przez $\|h\|_H$ dąży do zera.
Zatem $J$ jest różniczkowalny Frécheta, a jego pochodna ma postać
$DJ(u)v=\int_0^1u'(t)v'(t)\,\mathrm dt-
\int_0^1 f(t)v(t)\,\mathrm dt$.
Twierdzenie~\ref{thm:convex-derivative16} daje
równoważność $DJ(u)=0$ i globalnej minimalności.
Jeśli $f\equiv1$, funkcja $u(t)=t(1-t)/2$ spełnia
$u(0)=u(1)=0$ oraz $u''(t)=-1$. Ponieważ $u'\in L^2(0,1)$
i $\int_0^1u'(t)\,\mathrm dt=u(1)-u(0)=0$, mamy $u\in H$.
Dla absolutnie ciągłego $v\in H$ iloczyn $u'v$ jest
absolutnie ciągły i ma pochodną $u''v+u'v'$ prawie wszędzie.
Całkując i korzystając z $v(0)=v(1)=0$, otrzymujemy
$\int_0^1u'v'\,\mathrm dt=[u'v]_0^1-\int_0^1u''v\,\mathrm dt
=\int_0^1v\,\mathrm dt$.
Przez jednoznaczność jest to szukane rozwiązanie.
\end{proof}

\subsection{Wypukłe maksima i twierdzenie Kreina--Milmana}

Twierdzenie Kreina--Milmana (Twierdzenie~\ref{thm:krein-milman})
i zasada maksimum Bauera (Wniosek~\ref{cor:bauer9}) orzekają,
że niepusty zwarty wypukły podzbiór przestrzeni lokalnie wypukłej
Hausdorffa jest domkniętą otoczką wypukłą swoich punktów
ekstremalnych, a każda wypukła i półciągła z góry funkcja
rzeczywista na tym zbiorze osiąga maksimum w punkcie ekstremalnym. Ich rola w optymalizacji jest odmienna od metody bezpośredniej:
\emph{zwartość wypukłego zbioru} daje możliwość badania maksimum
funkcjonału wypukłego na punktach ekstremalnych,
natomiast \emph{koercywność i słaba półciągłość z dołu}
służą do znajdowania minimum.

\begin{stwierdzenie}[Maksimum funkcjonału afinicznego]
\label{prop:affine-extreme16}
Niech $E$ będzie lokalnie wypukłą przestrzenią liniową
topologiczną Hausdorffa nad $\K$, gdzie $\K=\R$ lub $\K=\C$,
$K\subseteq E$ niepustym, zwartym i wypukłym zbiorem,
a $A:K\to\R$ ciągłą funkcją afiniczną,
czyli $A((1-t)x+ty)=(1-t)A(x)+tA(y)$
dla wszystkich $x,y\in K$ i $t\in[0,1]$.
Wtedy istnieje punkt ekstremalny $e\in K$, w którym
$A$ osiąga maksimum, oraz (być może inny) punkt ekstremalny,
w którym osiąga minimum.
\end{stwierdzenie}
\begin{proof}
Funkcje $A$ i $-A$ są wypukłe i ciągłe,
więc zastosowanie zasady maksimum Bauera do obu daje tezę.
\end{proof}

\begin{przyklad}[Miary prawdopodobieństwa i ich punkty ekstremalne]
\label{ex:prob-measures16}
Niech $K$ będzie niepustą zwartą przestrzenią metryczną,
a $\mathcal P(K)$ zbiorem borelowskich miar probabilistycznych na $K$.
Niech $C(K;\R)$ będzie przestrzenią rzeczywistych funkcji
ciągłych z normą supremum. Na mocy twierdzenia
Riesza--Markowa--Kakutaniego (por. rozdziały 20 i 21
skryptu \emph{Miara i całka}) identyfikujemy $\mathcal P(K)$
z dodatnimi funkcjonałami liniowymi na $C(K;\R)$,
przyjmującymi wartość $1$ na funkcji stałej $1$.
Wszystkie skończone miary borelowskie na zwartej przestrzeni
metrycznej są regularne, więc reprezentacja obejmuje każdą
miarę z $\mathcal P(K)$.

Dla takiego funkcjonału $L$ z nierówności
$-\|\varphi\|_\infty 1\le\varphi\le\|\varphi\|_\infty 1$
i dodatniości wynika $|L(\varphi)|\le\|\varphi\|_\infty$.
Ponieważ $L(1)=1$, mamy $\|L\|=1$.
W topologii słabej-* $\sigma(C(K;\R)',C(K;\R))$
warunek $L(1)=1$ oraz każdy warunek $L(\varphi)\ge0$
dla $\varphi\in C(K;\R)$, $\varphi\ge0$, określają zbiory
domknięte, gdyż ewaluacje $L\mapsto L(\varphi)$ są ciągłe.
Ich przecięcie jest domkniętym podzbiorem słabo-* zwartej
kuli jednostkowej (Wniosek~\ref{cor:alaoglu-norm}),
a więc jest zwarte. Jest także wypukłe, gdyż dodatniość
i normalizacja są zachowane przez kombinacje wypukłe.

Punktami ekstremalnymi $\mathcal P(K)$ są dokładnie
miary Diraca $\delta_x$, $x\in K$.
Dla $\delta_x=t\mu+(1-t)\nu$, $0<t<1$, dodatniość
miar i równość ich mas wymuszają
$\mu(K\setminus\{x\})=\nu(K\setminus\{x\})=0$,
więc $\mu=\nu=\delta_x$.
Jeśli natomiast $\mu$ nie jest miarą Diraca,
istnieje zbiór borelowski $B$ o $0<\mu(B)<1$.
Aby to wykazać, niech $S$ oznacza nośnik $\mu$, czyli zbiór
punktów, których każde otwarte otoczenie ma dodatnią miarę.
Ponieważ zwarta przestrzeń metryczna ma przeliczalną bazę,
$K\setminus S$ jest sumą przeliczalnie wielu elementów bazy
o mierze zero. Stąd $\mu(K\setminus S)=0$ i $\mu(S)=1$.
Gdyby $S=\{x\}$, mielibyśmy $\mu=\delta_x$.
Zatem $S$ zawiera dwa różne punkty $x,y$.
Wybieramy ich rozłączne otwarte otoczenia $B,U$.
Z definicji nośnika $\mu(B)>0$ oraz $\mu(U)>0$,
więc $\mu(B)\le1-\mu(U)<1$.

Dla każdego borelowskiego $A\subseteq K$ określamy
$\mu_B(A)=\mu(A\cap B)/\mu(B)$ oraz
$\mu_{K\setminus B}(A)=\mu(A\setminus B)/(1-\mu(B))$.
Są to miary probabilistyczne, przy czym
$\mu_B(B)=1$ i $\mu_{K\setminus B}(B)=0$, zatem są różne.
Dla każdego $A$ suma $\mu(A\cap B)+\mu(A\setminus B)$
jest równa $\mu(A)$, więc
$\mu=\mu(B)\,\mu_B+(1-\mu(B))\,\mu_{K\setminus B}$.
Jest to nietrywialna kombinacja wypukła, zatem $\mu$
nie jest punktem ekstremalnym.
Twierdzenie Kreina--Milmana stwierdza więc, że
$\mathcal P(K)$ jest słabo-* domkniętą otoczką wypukłą
rodziny $\{\delta_x:x\in K\}$.
W szczególności dla $\varphi\in C(K;\R)$
maksimum funkcjonału $\mu\mapsto\int_K\varphi(x)\,\mathrm d\mu(x)$
można osiągnąć na pewnej mierze Diraca: funkcjonał ten
jest afiniczny i słabo-* ciągły jako ewaluacja na $\varphi$,
więc stosuje się Stwierdzenie~\ref{prop:affine-extreme16}.
\end{przyklad}

\begin{uwaga}[Dalsza perspektywa]
Twierdzenie Kreina--Milmana jest rezultatem o \emph{zwartych
zbiorach wypukłych} w ogólnych przestrzeniach lokalnie wypukłych,
a nie jedynie o kulach jednostkowych przestrzeni Banacha.
Z kolei subróżniczki i operatory proksymalne prowadzą
do teorii operatorów monotonicznych i metod minimalizacji,
co wykracza poza podstawowy zakres obecnego rozdziału.
\end{uwaga}

\end{document}
